

%%%%%%%%%%%%%%%%%%%%%%%%%%%%%
\begin{exo}[PC]{Défaut sur une lame de verre}
On considère une lame de verre à faces parallèles, transparente d'indice $n$ d'épaisseur $e_0$, éclairée par une source monochromatique de longueur d'onde dans le vide $\lambda_0$, située à l'infini (éclairage parallèle) et placée dans l'air d'indice $1$. Cette lame présente localement un défaut où l'épaisseur devient $e>e_0$.
\begin{quest}
\item Faire un schéma puis déterminer le déphasage à l'infini entre un rayon qui ne passe pas par le défaut et un rayon qui passe par le défaut.
\item Représenter une surface d'onde avant la traversée de la lame et après la traversée de la lame, et préciser ses caractéristiques en fonction des variables du problème.
\end{quest}
\end{exo}
%%%%%%%%%%%%%%%%%%%%%%%%%%%%%%%%%%%%%%%%%%%%%%%%%%%%%%%%%%%%%%%%


%%%%%%%%%%%%%%%%%%%%%%%%%%%%%%%%%%%
\begin{exo}[PC][3][TD]{Principe de l'holographie}
  \Notions{\item Interférences \task Source polychromatique \task Cohérence temporelle}
  L'holographie a été inventée par D. \textsc{Gabor} en 1947. Cette technique consiste à imprimer sur un support la figure d'interférence entre une onde de référence et une onde diffusée par un objet, de sorte à reproduire une image en trois dimensions de l'objet lorsque le support est éclairée par la seule onde de référence. Pour illustrer le principe de l'holographie, nous étudions un objet ponctuel situé au point O en aval du support, et éclairé par une onde plane monochromatique de longueur d'onde $\lambda$ normale au support. 
  \begin{fig}[label]
    \centering
    \includegraphics[scale=1]{hologramme.pdf}
    %\capt{}
  \end{fig}
  \begin{quest}
    \item	Justifier qu'il est légitime de supposer que l'onde diffusée par l'objet en O est une onde sphérique divergente de centre O. Rappeler l'expression de la vibration lumineuse associée $\vib(M,t)$.
    \sol{
    L'onde émise par un objet ponctuelle doit avoir une géométrie sphérique. Cette onde est divergente puisque le point $\rm O$ est la source de l'onde. La vibration lumineuse s'écrit, à une origine des phases près,
    \eqbox{
    s(M,t) = \frac{A}{OM} \cos\left(\omega t - \frac{2\pi}{\lambda}\opt{OM} \right).
    } 
    }
    \item Déterminer le chemin optique $\opt{OM}$ de l'onde diffusée par l'objet en un point M du support holographique. Déterminer aussi le chemin optique $\opt{HM}$ d'un rayon de l'onde de référence passant par le point $\rm H$ appartenant au plan d'onde de $\rm O$ et arrivant en $\rm M$.
    \sol{
    Le milieu est de l'air, d'indice assimilé à $1$. 
    \begin{fig}[label]
      \centering
      \includegraphics[scale=1]{hologramme_sol}
      %\capt{}
    \end{fig}
    Le chemin optique est donné par la distance parcourue par la lumière dans le milieu, soit par le théorème de Pythagore,
    \eqbox{
      \opt{OM} = \sqrt{d^2+x^2},
    }
    et simplement,
    \eqbox{
      \opt{HM} = d.
    }
    } 
    \item En déduire l'expression de la différence de marche $\delta(M)$ dans l'approximation paraxiale $x\ll d$, puis l'expression de l'intensité lumineuse $I(M)$ en fonction de l'abscisse $x$ du point M.
      \sol{
    La différence de marche est définie par la différence de phase entre les deux ondes en M, soit  ici 
    $$
    \delta(M) = \opt{OM} - \opt{HM} = \sqrt{d^2+x^2} - d = d\left[ (1+\veps)^{1/2}-1 \right] \simeq d \frac{\veps}{2},
    $$
    avec $\veps = x^2/d^2 \ll 1$ dans l'approximation paraxiale. On en déduit,
    \eqbox{
    \delta(M) = \frac{x^2}{2d}.
    }
    En supposant les deux ondes cohérentes (provenant d'une même source), l'intensité lumineuse s'obtient par la formule de Fresnel des interférences à deux ondes. En supposant de plus que les deux ondes sont de même intensité $I_0$, ce qui est possible dans l'approximation paraxiale si $A = I_0 d$, on a 
    $$
      I(M) = I_0\left[ 1 + \cos\left( \frac{2\pi}{\lambda} \delta(M) \right) \right],
    $$
    soit 
    \eqbox{
      I(x) = 2I_0\left[ 1 + \cos\left( \frac{\pi x^2}{\lambda d} \right) \right].
    }
    }
  \end{quest}
  Le support holographique est constitué d'une émulsion photosensible dont le coefficient de transmission en amplitude $t \in [0,1]$ est une fonction affine de l'intensité lumineuse reçue, soit 
  \begin{equation}
    t(M) = 1 - \frac{I(M)}{I_{\rm max }}.
  \end{equation}
   On retire maintenant l'objet tout en éclairant le support par l'onde de référence, comme ci-dessous.
  \begin{fig}[label]
    \centering
    \includegraphics[scale=1]{hologramme_vision.pdf}
    %\capt{}
  \end{fig}
  \begin{quest}
    \item	Justifier que le support holographique est assimilable à un réseau en transmission de pas variable~$a(x)$, dont on donnera l'expression.
    \sol{
    Traçons le graphe de la fonction $t(x)$. Sachant que $I_{\rm max} = 4I_0$, on trouve 
    $$
    t(x) = 1 - \frac{I(x)}{4I_0} = \frac{1}{2}\left[ 1 - \cos\left( \frac{\pi x^2}{\lambda d} \right) \right].
    $$
    La fonction $t(x)$ est représentée ci-dessous. On peut approximer $t(x)$ comme une successions de créneaux, ce qui correspond physiquement à une alternance de fentes (transmittance de $1$) et caches (transmittance de $0$), i.e. un réseau de diffraction.
    \begin{fig}[label]
      \centering
      \includegraphics[width=0.7\columnwidth]{hologramme_t.pdf}
      %\capt{}
    \end{fig}
    Le pas du réseau est donné par la distance entre deux fentes adjacentes, soit 
    $$
    I(a+x) = I(x) \Rar \delta(a+x) = \delta(x) + \lambda \Rar \frac{(a+x)^2}{2d} = \frac{x^2}{2d} + \lambda \Rar \frac{ax}{d} + \frac{a^2}{2d} = \lambda.
    $$
    Dans l'approximation $a(x) \ll x$, ce qui est valable en dehors de l'axe optique, on a
    \eqbox{
      a(x) = \frac{\lambda d}{x}.
    }
    }
  
    \item D'après le formule des réseaux, déterminer l'angle $\theta(x)$ donnant la direction des interférences constructives d'ordre $1$. On admettra que les ordres suivants sont supprimés à cause de la dépendance sinusoïdale de la fonction $t(M)$, et que l'ordre $0$ est supprimé en inclinant l'onde de référence par rapport au support d'un angle adéquat.
    \sol{
    La formule des réseaux donne la direction $\theta(x)$ des interférences constructives d'ordre $p$. Pour une incidence normale, on a
    $$
    \sin (\theta(x)) = \frac{p\lambda}{a(x)}.
    $$
    Pour l'ordre $p=1$, on trouve d'après l'expression du pas du réseau,
    \eqbox{
    \sin( \theta(x)) = \frac{x}{d}.
    } 
    }
    \item D'où semblent provenir les rayons lumineux émergents du support holographique dans la direction $\theta(x)$ ? En déduire que l'observateur voit une image virtuelle tridimensionnelle de l'objet en O.
    \sol{
    Dans l'approximation paraxiale, on trouve 
    $$
 \sin(\theta(x)) = \frac{x}{d} \simeq \tan(\theta(x)).
    $$
    Donc $\theta(x)$ correspond également à l'angle sous lequel un objet serait vu derrière le support holographique, à une distance $d$ et à une hauteur $x$. Cela correspond à la position précédente du point objet $\rm O$, qui rappelons-le, n'est plus présent dans ce second dispositif ! Nous venons de réaliser un hologramme.
    }
  \end{quest}
  \end{exo}
  %%%%%%%%%%%%%%%%%%%%%%%%%%%%%%%%%%%
  

%%%%%%%%%%%%%%%%%%%%%%%%%%%%%%%%%%%%%%%%%%%%%%%%%%%%%%%%%%%%%%%%%%%%%%%
\begin{exo}[PC]{Mesure de longueur d'onde avec les trous de Young}{1}
On cherche à mesurer la longueur d'onde d'un laser inconnu à l'aide du dispositif des trous de Young. La distance entre les trous est égale à $\rho$, et celle avec l'écran de projection est $Z=\SI{1,5}{m}$.
\begin{enumerate}
\item Dans un premier temps on dispose d'un laser HeNe rouge, pour lequel on connaît $\lambda_1=\SI{632,8}{nm}$. Sur l'écran, on mesure un écart de \SI{9,5}{mm} pour 6 interfranges. En déduire $\rho$ et estimer l'incertitude sur la mesure.
\item Sans modifier le reste du montage, on remplace le laser étalon par le laser dont on cherche à évaluer la longueur d'onde $\lambda_2$. On mesure cette fois un écart de \SI{8,0}{mm} toujours pour 6 interfranges. En déduire $\lambda_2$, et estimer l'incertitude sur la mesure.
\item On éclaire maintenant le dispositif avec une lumière blanche assimilée à une source ponctuelle. On observe sur l'écran des franges irisées. Expliquer.
\end{enumerate}
\end{exo}
%%%%%%%%%%%%%%%%%%%%%%%%%%%%%%%%%%%%%%%%%%%%%%%%%%%%%%%%%%%%%%%%%%%%%%%

%%%%%%%%%%%%%%%%%%%%%%%%%%%%%%%%%%%%%%%%%%%%%%%%%%%%%%%%%%%%%%%%%%%%%%%
\begin{exo}[PC][1]{Miroir de Lloyd}
Le dispositif de Lloyd permet d'obtenir des interférences à deux ondes. Il consiste en un miroir plan (par exemple horizontal) et un écran (par exemple vertical), éclairés par une source $S$ supposée ponctuelle et monochromatique de longueur d'onde $\lambda$ placée très proche du miroir. La source est à une hauteur $h$ au-dessus du miroir. Les rayons issus directement de la source et ceux issus de la réflexion sur le miroir interfèrent sur l'écran. On rappelle que la réflexion sur le miroir entraîne un déphasage supplémentaire de $\pi$ de l'onde réfléchie.
\begin{quest}
\item Représenter le champ d'interférences par une construction géométrique.
\item Montrer par une construction géométrique que le dispositif est équivalent à des trous de Young. On pourra faire intervenir l'image $S'$ de la source par le miroir.
\item Déterminer au point $M$ la différence de marche $\delta$, l'ordre d'interférences $p$ et l'intensité $I(x)$. En déduire l'interfrange $i$.
\item On éclaire maintenant le miroir avec une lampe à vapeur de sodium qui émet un doublet de longueur d'ondes $\lambda_1=\SI{589,0}{nm}$ et $\lambda_2=\SI{589,6}{nm}$. Déterminer l'expression de l'intensité lumineuse en fonction de la longueur d'onde moyenne $\lambda_m = \frac{1}{2}( \lambda_1 + \lambda_2)$ et de l'écart $\Delta \lambda = \lambda_2 - \lambda_1 \ll \lambda_m$. Déterminer le contraste et décrire la figure d'interférences.
% \item On rapproche la source du miroir d'une distance $\Delta h$ et on mesure un nouvel interfrange $i'=1,5 i$. Décrire la nouvelle figure d'interférences et exprimer la longueur d'onde $\lambda$ en supposant $\Delta h$ connu.
% \item On remplace la source ponctuelle par une fente. Que devient la figure d'interférences ?
\end{quest}
\end{exo}
%%%%%%%%%%%%%%%%%%%%%%%%%%%%%%%%%%%%%%%%%%%%%%%%%%%%%%%%%%%%%%%%%%%%%%%

%%%%%%%%%%%%%%%%%%%%%%%%%%%%%%%%%%%%%%%%%%%%%%%%%%%%%%%%%%%%%% 
\begin{exo}[PC][1][colle]{Miroirs de Fresnel} 
  On considère deux miroirs placé en coin et orienté d'un petit angle $\alpha$ l'un par rapport à l'autre. Une source ponctuelle $S$ placé à une distance $d$ du coin commun aux deux miroirs émet un faisceau de lumière monochromatique de longueur d'onde $\lambda$. Le milieu est de l'air d'indice de réfraction $n=1$.
\begin{fig}[]
  \centering
  \includegraphics[width=0.35\columnwidth]{miroirs_de_fresnel.pdf}
  % \capt{}
\end{fig}
  \begin{quest} 
    \item Représenter le trajet des rayons lumineux et construire géométriquement le champ d'interférences. En déduire la position des sources secondaires $S_1$ et $S_2$ ainsi que leur écartement~$a$. 
    \item Par analogie avec le dispositif des trous de Young, déterminer l'expression de l'intensité $I(x)$ sur un écran placé à une distance $D$ des miroirs. 
  \end{quest} 
\end{exo} %%%%%%%%%%%%%%%%%%%%%%%%%%%%%%%%%%%%%%%%%%%%%%%%%%%%%%%%%%%%%%

%%%%%%%%%%%%%%%%%%%%%%%%%%%%%%%%%%%%%%%%%%%%%%%%%%%%%%%%%%%%%%%%%%%%%%%
\begin{exo}[PC][3][TD]{Interférométrie stellaire}
  \Notions{\item Trous de Young \task Cohérence spatiale}
On observe une étoile à l'infini à travers une lunette astronomique. Un système de deux trous de Young d’écartement $a$ réglable est placé contre l'objectif, de distance focale $f>0$. On observe la figure d’interférences sur un écran dans le plan focal image de l’objectif. Un dispositif non représenté filtre la lumière incidente de l'étoile à la longueur d'onde de la raie jaune du sodium $\lambda=\SI{589.0}{nm}$, d'intensité $I_0$. L'ensemble du dispositif interférométrique est placé dans l'air, d'indice assimilé à $1$.
%\begin{figure}[h!]
%\centering
\begin{center}
\includegraphics[width=0.8\textwidth]{telescope.pdf}
\end{center}
%\end{figure}
\begin{quest}
\item Qu'observerait-on si l'étoile était ponctuelle et située à l’infini sur l’axe optique ? Donner l’expression de l’intensité lumineuse $I(x)$ sur l’écran. Préciser les caractéristiques de la figure d'interférences observée (forme des franges, espacement). Que se passe-t-il si l’écartement $a$ entre les trous augmente ?
\sol{
Il s'agit des trous de Young \guill{classique} en montage de Fraunhofer. Voir cours. L'intensité a pour expression 
\eqbox{
  I(x) = 2I_0\left[ 1+\cos(2\pi ax/\lambda f)\right].
  }
  Penser à représenter le graphe de $I(x)$ et à caractériser la figure d'interférences : franges rectilignes d'interfrange $i=\lambda f/a$ qui diminue si $a$ augmente.
}
\end{quest}
L'étoile n'est pas ponctuelle en réalité et on cherche à évaluer son angle apparent, c’est-à-dire l’angle $\alpha$ sous laquelle elle est vue à travers l’objectif. On considère pour commencer une source ponctuelle fictive située hors de l'axe optique.
\begin{quest}
  \item Que vaut l’intensité $I_\beta(x)$ issue d'une source ponctuelle à l’infini et vue sous un angle $\beta$ par rapport à l'axe optique ?
  \sol{
    Deux méthodes. Soit l'on détermine la nouvelle différence de marche $\delta$ à partir du chemin optique et d'un schéma similaire au montage de Fraunhofer. Soit l'on utilise la formule du cours sur la translation d'une source dans une direction parallèle aux trous. On sait que la nouvelle différence de marche vaut 
    $$
\delta' = \delta + \delta_t = \frac{ax}{f} - \frac{ab}{d}
    $$
    ou $d$ est la distance à la source et $b=d\tan\beta \simeq d \beta$ la hauteur. On en déduit 
    \eqbox{
    \delta' = \frac{ax}{f} - \frac{ab}{d} = a\left(\frac{x}{f} -  \beta\right).
    }
    
  }
\item  En modélisant l'étoile par une assemblée de sources ponctuelles incohérentes distribuées sur une largeur angulaire $\alpha$ et émettant la même intensité, montrer que l'intensité sur l’écran se met sous la forme
$$
I(x) = 2 I_0 \left[ 1 + \sinc{}\left(\pi \frac{a\alpha}{\lambda}\right) \cos\left(2\pi \frac{a x}{\lambda f}\right)\right],
$$ 
où $\sinc{} : x \mapsto \frac{\sin(x)}{x}$ est la fonction sinus cardinal.
Représenter le graphe de la fonction de contraste $\Co(a)$ et commenter. Définir en particulier une longueur de cohérence spatiale $\ls$.
\sol{
Imaginons une source élémentaire située à la distance $d$ des trous et à la hauteur $b$ au-dessus de l'axe optique. D'après le cours, la figure d'interférences est translatée sur l'écran d'une quantité $\Delta x = -\frac{b}{d} f$. Si la source est à l'infini, on remplace $-\frac{b}{d}$ par l'angle apparent $\beta$ (défini positif dans le sens trigo). Une source élémentaire de taille angulaire $\d \beta$ contribue à l'intensité d'une quantité
$$
\d I = A \left[ 1+\cos(2\pi a/\lambda f\times(x+\beta f))\right] \d \beta,
$$
avec $A$ une constante (intensité angulaire). Par changement de variable $u=2\pi a/\lambda f\times(x+\beta f)$, l'intégration donne
\begin{align*}
  I & = \int_{-\alpha/2}^{\alpha/2} A \left[ 1+\cos(2\pi a/\lambda f\times(x+\beta f))\right] \d \beta \\ & =  A\times\frac{\lambda}{2\pi a} \int_{2\pi a/\lambda f\times(x- \alpha f'/2)}^{2\pi a/\lambda f\times(x+\alpha f/2)} \left[ 1+\cos(u)\right] \d u \\
  & = A\times\frac{\lambda}{2\pi a} \left[ 2\pi a \alpha f/\lambda D +\sin(2\pi a/\lambda f\times(x+\alpha f/2)) - \sin(2\pi a/\lambda f\times(x-\alpha f/2))\right].
\end{align*}
Or $\sin(a)-\sin(b)=2\sin((a-b)/2)\cos((a+b)/2)$ et donc
$$
I = A\alpha \left[ 1 + \frac{\lambda}{\pi a \alpha}  \sin(\pi a\alpha/\lambda) \cos(2\pi ax/\lambda f)\right]
$$
On peut mettre sous la forme
\eqbox{
  I = 2I_0\left[1+\mc{C}(a) \cos\pa{2\pi \frac{x}{i}} \right],
}
où le contraste vaut 
\eqbox{
\Co(a) = \sinc \left(\pi \frac{a \alpha}{\lambda}\right).
}
Le graphe de $\Co(a)$ est celui de la figure ci-dessous. On remarque que le contraste diminue quand $a$ augmente, et s'annule pour $a=\ls$.
\begin{fig}[label]
  \centering
  \includegraphics[width=0.7\columnwidth]{contraste_spatial.pdf}
  %\capt{}
\end{fig}
}
\item Que se passe-t-il si l'écartement $a$ augmente à partir de zéro à $\alpha$ fixé ?
L'observation de l’étoile Betelgeuse par A. \textsc{Michelson} et F. \textsc{Pearse} en 1921 montre la présence de franges dont le contraste diminue quand on éloigne les deux trous l'un de l'autre, et s'annule une première fois pour $a=\SI{16.8}{cm}$. En déduire le diamètre apparent $\alpha$ de l'étoile, en radian puis en seconde d'arc. Déterminer son rayon $R$, sachant que sa distance à la Terre est $d\simeq\SI{200}{pc}$ (parsec), et où $1\,\text{pc}=\SI{3.1e16}{m}$. Commenter.
\sol{
  Le contraste diminue jusqu'à s'annuler puis changer de signe. De même si on augmente $a$.
Le contraste s'écrit aussi $\mc{C}(a) = \text{sinc}\left(  {\pi \frac{a}{l_s}}\right)$ où $l_s = \lambda/\alpha$ est la longueur de cohérence spatiale de la source. Comme $\mc{C}(a)$ s'annule pour $a=l_s =\SI{10}{cm}$, on a
\eqbox{
  \alpha = \frac{\lambda}{l_s} =
}
}
\end{quest}
\end{exo}
%%%%%%%%%%%%%%%%%%%%%%%%%%%%%%%%%%%%%%%%%%%%%%%%%%%%%%%%%%%%%%%%%%%%%%%

%%%%%%%%%%%%%%%%%%%%%%%%%%%%%%%%%%%%%%%%%%%%%%%%%%%%%%%%%%%%%%%%
\begin{exo}[PC]{Translation des franges}
On considère un dispositif de fentes de Young notées $F_1$ et $F_2$ distantes de $a=\SI{1}{mm}$, éclairées en incidence normale par une source $S$ ponctuelle, monochromatique de longueur d'onde $\lambda=\SI{633}{nm}$.

\begin{quest}
\item Qu'observe-t-on sur un écran placé à une distance $D=\SI{2}{m}$ ?
\item On suppose que l'on place une lame de mica transparente, à faces parallèles, d'épaisseur $e=\SI{50}{\micro\meter}$ et d'indice $n=1.5$ pour la radiation de longueur d'onde $\lambda$ utilisée, devant \textbf{une} seule des fentes. Qu'observe-t-on sur l'écran ? Expliquer comment est modifiée la figure d'interférences par rapport au cas précédent.
\end{quest}
\end{exo}
%%%%%%%%%%%%%%%%%%%%%%%%%%%%%%%%%%%%%%%%%%%%%%%%%%%%%%%%%%%%%%%%

%%%%%%%%%%%%%%%%%%%%%%%%%%%%%%%%%%%%%%%%%%%%%%%%%%%%%%%%%%%%%%%%
\begin{exo}[PC]{Fentes de Young éclairées en lumière blanche}
On éclaire deux fentes de Young, supposées infinies dans la direction $z$ et séparées d'une distance $b$ dans la direction $y$, avec une lampe quartz-iode assimilée à une source de lumière blanche. Les interférences sont observées à l'infini grâce à une lentille convergente de distance focale $f'$. On place à une distance $d$ des deux fentes une fente simple, également supposée infinie dans la direction $z$.
\begin{quest}
\item Justifier pourquoi on peut raisonner dans un plan $z=\cte$.
\item Représenter un schéma du montage avec la marche de deux rayons interférant en un point $M$ de l'écran.

On place maintenant un filtre interférentiel entre la lampe et la fente source. On supposera donc que le montage est éclairé par une lumière monochromatique de longueur d'onde $\lambda_0$.

\item En supposant la fente source infiniment fine et placée sur l'axe optique, établir l'expression de l'ordre d'interférences $p_0(y)$. En déduire la position des franges brillantes en fonction d'un entier $m$ et de l'interfrange $i$.
\item Pour obtenir une figure d'interférences de grande taille, comment doit-on choisir la focale $f'$ ? Expérimentalement, quel peut-en être l'inconvénient ? À quelle distance des bifentes doit-elle être placée ?

On suppose maintenant que la fente source infiniment fine est décalée d'une distance $a\ll d$ par rapport à l'axe optique.

\item Établir la nouvelle expression de l'ordre d'interférences $p_a(y)$. Comment la figure d'interférences est-elle modifiée ?

On élargit maintenant la fente source, qui devient de largeur $2a$ dans la direction $y$.

\item Expliquer pourquoi la figure d'interférences se brouille si $p_a(y)-p_0(y)\gg 1$.

\item Déterminer la largeur de cohérence spatiale $l_{\rm{s}}$ telle que $p_{l_{\rm{s}}}(y)-p_0(y)=\frac{1}{2}$. Commenter le résultat: la largeur $l_{\rm{s}}$ est-elle intrinsèque à la source ?

On enlève maintenant le filtre interférentiel, et on réduit la largeur de la fente source de sorte que $a\ll l_{\rm{s}}$. La figure d'interférences observée est représentée ci-dessous.

\begin{center}
\includegraphics[width=0.5\textwidth]{interfs_lumiere_blanche.png}
\end{center}

\item En s'appuyant sur les questions précédentes, interpréter qualitativement la figure. Pourquoi la frange centrale est-elle blanche ? Pourquoi obtient-on des irisations ?

\item Si l'on élargit la fente source de sorte que $a>l_{\rm{s}}$, qu'observe-t-on sur l'écran ?
\end{quest}
\end{exo}
%%%%%%%%%%%%%%%%%%%%%%%%%%%%%%%%%%%%%%%%%%%%%%%%%%%%%%%%%%%%%%%%

%%%%%%%%%%%%%%%%%%%%%%%%%%%%%%%%%%%%%%%%%%%%%%%%%%%%%%%%%%%%%%%%
\begin{exo}[PC][2][TD]{Biprisme de Fresnel}
  \Notions{\item Champ d'interférences \task Trous de Young}
Un biprisme est éclairé par une fente fine de centre $\rm S$, située dans le plan de symétrie du biprisme. La lumière est monochromatique de longueur d'onde $\lambda_0=\SI{589}{nm}$. La source $\rm S$ est à une distance $d=\SI{10}{cm}$ du biprisme très grande devant la taille du biprisme. On note $A=\SI{11.2e-3}{rad}$ l'angle du biprisme et $D=\SI{6}{m}$ la distance entre le biprisme et l'écran d'observation. Le verre est d'indice de réfraction $n=\SI{1.52}{}$.
%\begin{figure}[h!]
%\centering
\begin{center}
  % \includegraphics[width=0.5\textwidth]{biprisme.pdf}
  \includegraphics[width=0.3\columnwidth]{biprisme_angle.pdf}
\end{center}
%\end{figure}
\begin{quest}
\item En supposant tous les angles petits devant $1$, montrer que tout rayon arrivant sur la face d'entrée est dévié de $\alpha=(n-1)A$. Application numérique.
\sol{
Comme $d$ est très grande devant la taille du biprisme, l’angle d’incidence $i$ est petit. D’après Snell-Descartes, $i \simeq \sin i = n \sin r \simeq n r$ pour la réfraction en entrée, et  $ n i’ \simeq n \sin i’ = \sin r’ \simeq r’$ pour celle en sortie. Un peu de géométrie donne que $r + i’ = A/2$. La déviation angulaire vaut $\alpha = (r-i) + (r’-i’) = (n-1) r + (n-1)i’ = (n-1)(r+i’) = (n-1)A$, d'où le résultat,
\eqbox{
  \alpha=(n-1)A.
}
}
\item Représenter le trajet des rayons lumineux et construire géométriquement le champ d'interférences. En déduire la position des sources secondaires $\rm S_1$ et $\rm S_2$ ainsi que leur écartement~$a$. 
\sol{
  Au premier ordre en $\alpha$, les sources secondaires $S_1$ et $S_2$ doivent àtre situées dans le même plan que la source $S$, puisque leur positions sont symétriques par rapport à l'axe optique, et doivent coincider avec la source $S$ si $\alpha=0$. 
  \begin{fig}[]
    \centering
    \includegraphics[width=0.75\textwidth]{biprisme.pdf}
    % \capt{}
  \end{fig}
L'écartement vaut
\eqbox{
  a \simeq 2\alpha d = \SI{1.75}{mm}
}
}
\item Par analogie avec les trous de Young, déterminer le profil d'intensité sur l'écran $I(x)$ et décrire la figure d'interférences. Combien de franges peut-on espérer voir dans le champ d'interférences ? Comment les observe-t-on en pratique ? 
\sol{
Tout comme les trous de Young, mais la distance $D$ correspond ici à la distance entre les sources secondaires et l'écran.
}
% On prendra $A=\SI{1}{\degree}$, $d=\SI{10}{cm}$, $D=\SI{1}{m}$ et $n=\num{1.5}$.
\end{quest}
\end{exo}
%%%%%%%%%%%%%%%%%%%%%%%%%%%%%%%%%%%%%%%%%%%%%%%%%%%%%%%%%%%%%%%%


%%%%%%%%%%%%%%%%%%%%%%%%%%%%%%%%%%%
\begin{exo}[PC][2][TD]{Déviation minimale des rayons dispersés par un réseau}
  \Notions{\item Réseau de diffraction \task Déviation angulaire}
On considère un réseau de fentes éclairé par une onde plane monochromatique de longueur d'onde $\lambda$ faisant un angle $i$ avec la normale au réseau. On observe les rayons diffractés dans une direction repérée par l'angle de déviation $\alpha$, i.e. l'angle par rapport à la direction des rayons incidents. Ces rayons interfèrent entre eux à l'infini.
\begin{quest}
  \item Faire un schéma. Exprimer la condition d'interférence constructive entre les rayons diffractés. En déduire la loi donnant la déviation $\alpha$ correspondant aux pics principaux.
  \sol{
    Sous une incidence $i$, le rayon arrive en $\mathrm{F}$ en avance par rapport à celui arrivant en F'. Cependant, il prend ensuite du retard de sorte que le déphasage s'écrit
    Il y a interférences constructives entre tous les rayons, lorsque deux ondes consécutives sont déphasées d'un multiple de $2 \pi$, soit
    \eqbox{
      \sin (i+\alpha)-\sin i=p \frac{\lambda}{a},~p \in \mathbb{Z}.
    }
    La déviation angulaire $\alpha$ dépend non seulement de la lumière incidente, du réseau et de l'ordre d'interférence $p$, mais également des conditions d'incidence (via l'angle $i$). L'ordre $p=0$ correspond aux rayons non déviés ($\alpha=0$).
  }
  \item Montrer que pour un ordre $p$ fixé, lorsque l'on fait varier l'angle $i$ la déviation $\alpha$ passe par un extremum que l'on exprimera.
  \sol{
    Dérivons la relation par rapport à l'angle $i$, l'ordre $p$ étant fixé :
    $$
    \left(1+\frac{\mathrm{d} \alpha}{\mathrm{~d} i}\right) \cos (i+\alpha)-\cos i=0
    $$
    
    La déviation $\alpha(i)$ passe par un extremum quand
    $$
    \frac{\mathrm{d} \alpha}{\mathrm{~d} i}=0 \quad \text { soit } \quad \cos (i+\alpha)=\cos i
    $$
    c'est-à-dire (en excluant $i=0$ qui correspond à l'ordre 0) 
    \eqbox{
    \alpha_{\mathrm{m}}=-2 i.
    }
    %  En dérivant à nouveau, on trouve $D^{\prime \prime}(\alpha)=-2 \tan \alpha$ lorsque $D=-2 \alpha . D_{\mathrm{m}}$ est donc minimum quand $\alpha<0$ et maximum quand $\alpha>0$. Dans tous les cas, $\left|D_{\mathrm{m}}\right|$ est minimum.
    
    En substituant dans la loi des réseaux, il vient
    \eqbox{
      2 \sin \left( \frac{\alpha_{\mathrm{m}}}{2} \right) = p \frac{\lambda}{a}
    }
  }
  \item Quelle est l'orientation du réseau par rapport aux faisceaux incident et diffracté d'ordre $p$, lorsque la déviation est minimale?
  \sol{
    La déviation est minimale lorsque le réseau formue une bissectrice entre le faisceau incident et le faisceau diffracté.
  }
\end{quest}
\end{exo}
%%%%%%%%%%%%%%%%%%%%%%%%%%%%%%%%%%%

%%%%%%%%%%%%%%%%%%%%%%%%%%%%%%%%%%%%%%%%%
\begin{exo}[PC][2][TD]{Bilentille de Billet}
  \Notions{\item Champ d'interférences \task Trous de Young}
On considère une lentille convergente $\mc{L}$ de distance focale $f'=\SI{20}{cm}$. On place une source ponctuelle monochromatique $S$ de longueur d'onde $\lambda_0=\SI{500}{nm}$ sur l'axe de la lentille, à une distance $SO=\SI{60}{cm}$ de son centre. On suppose que l'indice de l'air vaut $n_{\rm air}\simeq 1$.
\begin{quest}
  \item Construire géométriquement l'image $\rm S'$ de $\rm S$ par la lentille. Que vaut le grandissement ?
  \sol{
    On utilise les règles standards de construction en supposant qu’on place un objet $AB$ transversal en $S$.
    \begin{center}
      \includegraphics{bilentille_1.pdf}
    \end{center}
    }
  \end{quest}
  La lentille est sciée en deux parties identiques qu'on éloigne l'une de l'autre d'une distance $e=\SI{2}{mm}$. On place un cache entre les deux demi-lentilles ainsi obtenues. On note $\rm O_1$ et $\rm O_2$ les centres optiques respectifs de demi-lentilles supérieure et inférieure.
  \begin{fig}[]
    \centering
    \includegraphics[width=0.7\textwidth]{bilentilles_de_billet_simple.pdf}
    % \capt{}
  \end{fig}
\begin{quest}
  \item Déterminer géométriquement les images de $\rm S$, notées $\rm S_1$ et $\rm S_2$, par les moitiés de lentille supérieure et inférieure, et calculer la distance $a=S_1S_2$. Construire géométriquement le champ d'interférences.
  \sol{
  \begin{center}
    \includegraphics{bilentille_2.pdf}
  \end{center}
  \begin{fig}[]
    \centering
    \includegraphics[width=0.8\textwidth]{bilentilles_de_billet.pdf}
    % \capt{}
  \end{fig}

  Thalès donne $\frac{S_2S_1}{O_2O_1} = \frac{SS'}{SO}$ avec $SO=\SI{60}{cm}$ et $O_2O_1=e$ et $SS'=\SI{90}{cm}$. On trouve \fbox{$a=S_1S_2=\SI{3}{mm}$}.
  }
  \item Par analogie avec les trous de Young, déterminer le profil d'intensité sur l'écran $I(x)$ et décrire la figure d'interférences. Combien de franges peut-on espérer voir dans le champ d'interférences ? Comment les observe-t-on en pratique ? On donne $D=\SI{1}{m}$ la distance entre les sources secondaires et l'écran.
  \sol{
  On produit deux sources secondaires placées symétriquement autour de l’axe optique. D’un point de vue optique, on vient donc de réaliser des trous d’Young. Tout ce qui a été fait en cours sur les trous d’Young s’applique ici. Le champ d'interférences sur un écran situé à une distance $D$ derrière les sources secondaires est de largeur $L=a\frac{D+OS'}{OS'}$. Le nombre de franges visibles vaut donc 
  $N=2L/i$, soit
  \eqbox{
  N = \frac{2a^2}{\lambda D}\left(1+\frac{D}{OS'}\right).  
  }
  }
\end{quest}
\end{exo}
%%%%%%%%%%%%%%%%%%%%%%%%%%%%%%%%%%%%%%%%%%%%%%%%

%%%%%%%%%%%%%%%%%%%%%%%%%%%%%%%%%%%%%%%%%%%%%%%%%%%%%%%%%%%%%%%%
\begin{exo}[PC][1]{Les trous de Young}
\Notions{\item Trous de Young \task Différence de marche \task Figure d'interférence}
Deux trous aux positions $\rm S_1$ et $\rm S_2$ sont éclairés par une source $\rm S$ de lumière monochromatique de longueur d'onde dans le vide $\lambda_0$. Les trous sont espacés d'une distance $a$ dans la direction $x$. Les trous produisent une figure d'interférence sur un écran situé à une distance $D$ derrière les trous dans la direction $z$ de l'axe optique. Le milieu est de l'air  d'indice de réfraction~$n_{\rm air}$.

%%% FIGURE %%%
%\begin{figure}[h!]
%\centering
\begin{center}
\includegraphics[scale=1]{experience_young_simple.pdf}
\end{center}
%\caption{Les trous d'Young $S_1$ et $S_2$.}
%\label{fig:}
%\end{figure}
%%%%%%%%%%

\begin{quest}
% \item Justifier que l'onde lumineuse qui éclaire les trous peut être considérée comme plane. Les trous $S_1$ et $S_2$ appartiennent-ils à la même surface d'onde.

% \sol{
% Source lointaine $\rightarrow$ onde plane. $S 1$ et $S_2$ appartiennent au même front d'onde.
% }

\item Soit $M(x,y)$ un point de l'écran. Par un développement limité, déterminer la différence de marche $\delta$ dans l'approximation paraxiale $D\gg a,x,y$. En déduire le déphasage $\Delta \vphi$ en fonction de $a$, $x$, $D$ et la longueur d'onde dans l'air $\lambda$, que l'on définira à partir de $\lambda_0$ et $n_{\rm air}$.

\sol{
  À un signe global près, on a
  \begin{equation*}
    \delta(M) = \opt{S M}_1-\opt{S M}_2 = \opt{S_2 M}-\opt{S_1 M} = n(S_2M-S_1M),
    \end{equation*}
    car $\opt{SS_1} = \opt{SS_2}$ par symétrie.
    Notons $O$ l'origine de l'axe $x$ sur l'écran, situé sur l'axe optique, et $H_1$ et $H_2$ les projetés orthogonaux de $S_1$ et $S_2$ sur l'écran. On a alors d'après le théorème de Pythagore

  \begin{equation*}
    \begin{cases}
        S_1 M^2 & = S_1H_1^2 +H_1M^2 = D^2 + (x + \frac{a}{2})^2 + y^2 \\ 
        S_2 M^2 & = S_2H_1^2 +H_2M^2 = D^2 + (x - \frac{a}{2})^2 + y^2
    \end{cases}
    % \end{split}
    \end{equation*}
  Pour obtenir une expression approchée simple de $\delta(M)$, on peut effectuer l'approximation paraxiale et supposer que l'écran est loin des trous et que l'on ne considère que la zone de l'écran proche de l'axe optique, c'est-à-dire que $D \gg a,x,y$.
   Au premier ordre, on a le développement limité
  \begin{equation}
  S_1 M = D\sqrt{1+\underbrace{\left(\frac{a}{2D}\right)^2+\dfrac{x^2 + y^2 - ax}{D^2}}_{\veps \ll 1}}\simeq D\left(1 +\frac{\veps}{2}\right) = D + \dfrac{a^2+4x^2 + 4y^2}{8D}-\frac{ax}{2D},
  \end{equation}
  On montre de la même manière que
  \begin{equation}
  S_2 M\simeq D+\dfrac{a^2+4x^2 + 4y^2}{8D}+\frac{ax}{2D}.
  \end{equation}
  La différence de marche vaut donc
  \eqbox{
    \label{eq:deltaM}
    \delta(M)\simeq\frac{nax}{D} \Rar \Delta \vphi = 2\pi\frac{\delta(M)}{\lambda_0} = \frac{2\pi ax}{\lambda D},
  }  
  avec $\lambda = \lambda_0/n$ la longueur d'onde dans le milieu.
  

}
\item Retrouver l'expression approchée de la différence de marche $\delta$ dans l'approximation paraxiale en considérant que les rayons provenant de $S_1$ et $S_2$ sont parallèles et forment un angle $\theta$ avec l'axe optique. On s'appuiera sur un schéma pour calculer les chemins optiques.
\sol{
La démonstration est celle utilisée dans le cours avec le montage de Fraunhofer.
}

\item Les sources secondaires $S_1$ et $S_2$ émettent une lumière de même intensité $I_0$. Quelle est l'expression de l'intensité $I$ sur l'écran en fonction et $I_0$ et de $x$ ? Représenter le graphe de cette fonction et caractériser la figure d'interférences. Déterminer l'interfrange $i$.

\sol{
   En utilisant la formule de Fresnel, on trouve
  \eqbox{
    I(x)=2I_{0}\left[1+\cos\left(2\pi \frac{a x}{\lambda D}\right)\right].
  }
  L'intensité sur l’écran ne dépend donc pas de $y$, ce qui est cohérent en raison de la symétrie du dispositif. La figure d’interférences est ainsi composé de \important{franges rectilignes} dans la direction $y$, alternativement sombres et brillantes. Ces fentes sont régulièrement espacées et séparées d'une distance $i$ appelée \imp{interfrange}, d'expression 
  \eqbox{
    i=\frac{\lambda D}{a}.
  }
}
\end{quest}
\end{exo}
%%%%%%%%%%%%%%%%%%%%%%%%%%%%%%%%%%%%%%%%%%%%%%%%%%%%%%%%%%%%%%%%%%%%%%%%%



%%%%%%%%%%%%%%%%%%%%%%%%%%%%%%%%%%%%%%%%%%%%%%%%%%%%%%%%%%%%%%%%
\begin{exo}[PCSI][1]{Expériences des trous d'Young}

Deux trous $S_1$ et $S_2$ sont éclairés par une source $S$ de lumière monochromatique de longueur d'onde $\lambda$, supposée lointaine. Les trous sont espacés d'une distance $a$ dans la direction $x$. Les trous produisent une figure d'interférence sur un écran situé à une distance $D$. Le milieu est de l'air d'indice de réfraction~$n_{\rm air}=1$.

%%% FIGURE %%%
%\begin{figure}[h!]
%\centering
\begin{center}
\includegraphics[width=0.9\columnwidth]{experience_young.pdf}
\end{center}
%\caption{Les trous d'Young $S_1$ et $S_2$.}
%\label{fig:}
%\end{figure}
%%%%%%%%%%

\begin{quest}
% \item Rappele
\item Soit $M(x)$ un point de l'écran sur l’axe $(Ox)$. Par un développement limité, déterminer la différence de marche $\delta$ dans l'hypothèse $D\gg a$. En déduire le déphasage $\Delta \phi$.

\sol{
$\delta \simeq n (S1M-S2M) = ax/D$ et $\Delta \phi = 2\pi\delta/\lambda$.
}


\item Les sources secondaires $S_1$ et $S_2$ émettent une lumière de même intensité $I_0$. Quel est l'expression de l'éclairement $I$ sur l'écran en fonction et $I_0$ et $\Delta \phi$, puis de $x$ ? En déduire l'interfrange $i$.

\sol{
$I = I_0\left[1+\cos(2\pi a x/\lambda D)\right]$ et $i \simeq \lambda D/a$.
}

\item Soit $S(M,t)=S_m \cos(\omega t + \varphi(M))$ l’expression générale de l’amplitude lumineuse sur l’écran. On suppose maintenant que la source lumineuses émet un doublet de raies aux pulsations $\omega_1$ et $\omega_2$. Montrer l’existence d’un phénomène de battement. À quoi ressemble alors la figure d’interférence ?
\end{quest}
\end{exo}
%%%%%%%%%%%%%%%%%%%%%%%%%%%%%%%%%%%%%%%%%%%%%%%%%%%%%%%%%%%%%%%%%%%%%%%%%



%%%%%%%%%%%%%%%%%%%%%%%%%%%%%%%%%%%%%%%%%%%%%%%%%%%%%%%%%%%%%%%%%%
\begin{exo}[PC][1][colle]{Trous de Young et lame de verre}
  \Notions{\item Trous de Young \task Lame de verre}
Deux trous $S_1$ et $S_2$ sont éclairés par une source ponctuelle $S$ émettant lumière monochromatique de longueur d'onde~$\lambda$. Les trous $S_1$ et $S_2$ sont espacés d'une distance $a$ dans la direction $x$. On observe la figure d'interférences formée sur un écran situé à une distance $D$ dans la direction transverse $z$ derrière les trous. Le milieu est de l'air d'indice de réfraction~$n_{\rm air}=1$.

\begin{quest}
  \item	 Soit $M(x)$ un point de l'écran placé à une abscisse $x$ sur l’axe $(Ox)$. Par un développement limité, déterminer la différence de marche $\delta$ dans l'hypothèse paraxiale $D\gg a,x$. Montrer que le déphasage $\Delta \vphi$ entre les sources secondaires $S_1$ et $S_2$ s'écrit 
  \begin{equation}
    \Delta \vphi = \frac{2\pi a x}{\lambda D}.
  \end{equation}

  \sol{
  $\delta \simeq n (S_1 M-S_2  M) = ax/D$ et $\Delta \phi = 2\pi\delta/\lambda$.
  }
  
  
  \item On suppose que la diffraction par les trous est isotrope, de sorte que les sources secondaires $S_1$ et $S_2$ émettent une lumière uniforme de même intensité $I_0$ sur l'écran. Quel est l'expression de l'intensité résultante $I$ sur l'écran en fonction et $I_0$ et $\Delta \vphi$, puis de $x$, $a$ et $\lambda$ et $D$ ? En déduire l'interfrange $i$.
  
  \sol{
  $I = I_0\left[1+\cos(2\pi a x/\lambda D)\right]$ et $i \simeq \lambda D/a$.
  }
\end{quest}
On ajoute derrière $S_2$ une lame de verre d'indice de réfraction $n >1$ et d'épaisseur $e$.
  \begin{quest}
  \item Déterminer la nouvelle expression de la différence de marche $\delta$. Comment est affectée la figure d'interférence ? En particulier, que vaut l'interfrange ? Établir l'expression de la position $x_0$ de la frange centrale, d'ordre d'interférences nulle.
  
  \sol{
    La lame de verre introduit un chemin optique supplémentaire $(n-1)e$ dans le bras $2$. On en déduit 
    \eqbox{
      \delta = \frac{ax}{D} - (n-1)e.
    } 
  La différence de marche supplémentaire introduite par la lame de verre est indépendante de $x$. L'interfrange reste inchangé. Par contre la figure est translatée. La frange centrale est située maintenant en une position $x_0$ telle que $\delta = \frac{a}{D}(x-x_0)$, soit
  \eqbox{
  x_0 = (n-1)\frac{eD}{a}.
  }
  }
  
  \item On utilise comme source $S$ un laser de longueur d'onde $\lambda = \SI{632}{nm}$.	En ajoutant une lame de verre d'épaisseur $e=\SI{100}{\micro m}$, on observe que la frange centrale s'est décalée de $N=94$ franges. De quel type de verre s'agit-il (on s'appuiera sur le tableau suivant) ? La mesure dépend-elle de $a$ et $D$ ?
  
  \sol{
    La position de la frange centrale vaut ici $x_0 = N i = N \frac{\lambda D}{a}$. On en déduit
    \eqbox{
      n = 1+ \frac{N\lambda}{e} = \SI{1.59}{}.
    }
    Il s'agit d'un verre \emph{flint}.
  }

\end{quest}

\begin{center}
\includegraphics[width=0.45\textwidth]{index1.png}
\includegraphics[width=0.45\textwidth]{index2.png}
\end{center}

\end{exo}
%%%%%%%%%%%%%%%%%%%%%%%%%%%%%%%%%%%%%%%%%%%%%%%%%%%%%%%%%%%%%%%%%%%%%%%%%

%%%%%%%%%%%%%%%%%%%%%%%%%%%%%%%%%%%%%%%%%%%%%%%%%%%%%%%%%%%%%%%%%%
\begin{exo}[PC][1][TD]{Mesure de l'indice d'une lame de verre}
\Notions{\item Trous de Young \task Différence de marche \task Lame de verre}
On part du montage des trous de Young de l'exercice précédent. On suppose que l'air est d'indice de réfraction $n_{\rm air} \simeq 1$.  On ajoute derrière $\rm S_2$ une lame de verre d'indice de réfraction $n$ et d'épaisseur $e$.

\begin{fig}[label]
  \centering
  \includegraphics[scale=1]{experience_young_lame.pdf}
  %\capt{}
\end{fig}
\begin{quest}
\item Déterminer la nouvelle expression de la différence de marche $\delta$. Comment est affectée la figure d'interférence ? En particulier, que vaut l'interfrange ? Déterminer la position $x_0$ de la frange centrale.

\sol{
  La lame de verre introduit un chemin optique supplémentaire $(n-1)e$ dans le bras $2$. On en déduit 
  \eqbox{
    \delta = \frac{ax}{D} - (n-1)e.
  } 
La différence de marche supplémentaire introduite par la lame de verre est indépendante de $x$. L'interfrange reste inchangé. Par contre la figure est translatée. La frange centrale est située maintenant en une position $x_0$ telle que $\delta = \frac{a}{D}(x-x_0)$, soit
\eqbox{
x_0 = (n-1)\frac{eD}{a}.
}
}

\item On utilise comme source $\rm S$ un laser He-Ne de longueur d'onde $\lambda_0 = \SI{632}{nm}$.	En ajoutant une lame de verre d'épaisseur $e=\SI{100}{\micro m}$, on observe que la frange centrale s'est décalée de $N=94$ franges. De quel type de verre s'agit-il (on s'appuiera sur le tableau suivant) ? La mesure dépend-elle de $a$ et $D$ ?

\sol{
  La position de la frange centrale vaut ici $x_0 = N i = N \frac{\lambda D}{a}$. On en déduit
  \eqbox{
    n = 1+ \frac{N\lambda}{e} = \SI{1.59}{}.
  }
  Il s'agit d'un verre \emph{flint}.
}
\end{quest}
\begin{center}
\includegraphics[width=0.45\textwidth]{index1.png}
\includegraphics[width=0.45\textwidth]{index2.png}
\end{center}
\end{exo}
%%%%%%%%%%%%%%%%%%%%%%%%%%%%%%%%%%%%%%%%%%%%%%%%%%%%%%%%%%%%%%%%%%%%%%%%%

%%%%%%%%%%%%%%%%%%%%%%%%%%%%%%%%%%%%%%%%%%%%%%%
\begin{exo}[MPSI]{Les fentes d'Young}

Deux fentes $S_1$ et $S_2$ sont éclairés par une source $S$ de lumière monochromatique de longueur d'onde $\lambda$, supposée lointaine. Les fentes sont orientées selon $y$, de largeur $b$, et espacées d'une distance $a\gg b$ dans la direction $x$. Les fentes produisent une figure de diffraction et d'interférences sur un écran situé à une distance $D$. Le milieu est d'indice de réfraction~$n=1$.

%%% FIGURE %%%
%\begin{figure}[h!]
%\centering
\begin{center}
\includegraphics[scale=1]{experience_young.pdf}
\end{center}
%\caption{Les fentes d'Young $S_1$ et $S_2$.}
%\label{fig:}
%\end{figure}
%%%%%%%%%%

\begin{quest}
\item En comparant les distances $SS_1$ et $SS_2$ sur le schéma, justifier que l'onde lumineuse a même phase en $S_1$ et $S_2$. De quelle géométrie est l'onde lumineuse produite par $S$ ? Justifier que l'onde lumineuse qui éclaire les fentes peut malgré tout être considérée comme plane.

\sol{
Source lointaine $\rightarrow$ onde plane. $S 1$ et $S_2$ appartiennent au même front d'onde.
}


\item Diffraction. Dans l'hypothèse $D\gg a$, détermimer la taille $\ell$ sur l'écran de la figure de diffraction produite par chaque fente.

\sol{
$\ell \simeq \theta D \simeq \lambda D /b$.
}



\item Interférence.


\begin{enumerate}[label=\alph{*})]
\item Soit $M(x,y)$ un point de l'écran. Justifier que la figure d'interférence est invariante par translation selon $y$. Déterminer la différence de marche $\delta$ dans l'hypothèse $D\gg a$. En déduire le déphasage $\Delta \phi$.

\sol{
$\delta \simeq n (S1M-S2M) = ax/D$ et $\Delta \phi = 2\pi\delta/\lambda$.
}

\item Proposer une expression pour la valeur des vibrations lumineuses $s_1(x,t)$ et $s_2(x,t)$ sur l'écran, supposées sinusoïdales, issues de $S_1$ et $S_2$ respectivement. On notera $A$ l'amplitude (identique) de chacune des ondes. Calculer l'amplitude totale $s=s_1+s_2$.


\sol{
$s(x,t) = A\cos(\phi_0 - \omega t) + A\cos(\phi_0-\omega t +\Delta \phi )$ avec $\omega=2\pi c/\lambda$.
$$s(x,t) = 2A\cos\left(\phi_0 - \omega t + \dfrac{1}{2}\Delta \phi\right)\cos\left(\dfrac{1}{2}\Delta \phi\right)$$
}
\item L'intensité lumineuse $I$ s'obtient par la moyenne temporelle de la vibration lumineuse au carré, soit $I=\langle s^2 \rangle$. On donne $\langle \cos(a-\omega t)^2\rangle = \frac{1}{2}$. Quelle est l'expression de l'intensité lumineuse sur l'écran en fonction et $A$ et $\Delta \phi$, puis de $x$ ? Quel est le lieu des interférences destructives ? constructives ?  En déduire l'interfrange $i$.

\sol{
$I = A^2\cos\left(\dfrac{1}{2}\Delta \phi\right)^2 = \dfrac{A^2}{2}\left[1+\cos(\Delta\phi)\right] = \dfrac{A^2}{2}\left[1+\cos(2\pi a x/\lambda D)\right]$. $i \simeq \lambda D/a$.
}


\end{enumerate}

\item Figure complète. Relier l'intensité lumineuse $I_0$ produite par une des fentes à la grandeur $A$. Justifier qu'à cause de la diffraction, $I_0(x)$ dépend en réalité de $x$ et décroît sur une distance caractéristique $\ell$. En déduire l'allure du profil d'intensité sur l'écran. On fera apparaître les distances caractéristiques $i$ et $\ell$.

\sol{
$I_0 = \dfrac{A^2}{2}$ donc $I(x) = I_0(x)\left[1+\cos(2\pi a x/\lambda D)\right]$.
}

\end{quest}
\end{exo}
%%%%%%%%%%%%%%%%%%%%%%%%%%%%%%%%%%%%%%%%%%%%%


%%%%%%%%%%%%%%%%%%%%%%%%%%%%%%%%%%%
\begin{exo}[PC][1][devoir]{Interférences lumineuses : dispositif des trous d'Young}

  \Annale[Physique TSI]{CCINP}{2007}
% \section*{Données mathématiques :}
% Opérateurs mathématiques en coordonnées cartésiennes:

% Divergence :

% $$
% \operatorname{div} \vec{A}=\frac{\partial A_{x}}{\partial x}+\frac{\partial A_{y}}{\partial y}+\frac{\partial A_{z}}{\partial z}
% $$

% Rotationnel :

% $$
% \overrightarrow{\operatorname{rot}} \vec{A}=\left(\frac{\partial A_{z}}{\partial y}-\frac{\partial A_{y}}{\partial z}\right) \vec{e}_{x}+\left(\frac{\partial A_{x}}{\partial z}-\frac{\partial A_{z}}{\partial x}\right) \vec{e}_{y}+\left(\frac{\partial A_{y}}{\partial x}-\frac{\partial A_{x}}{\partial y}\right) \vec{e}_{z}
% $$

% Laplacien d'un champ scalaire : $\quad \Delta U=\frac{\partial^{2} U}{\partial x^{2}}+\frac{\partial^{2} U}{\partial y^{2}}+\frac{\partial^{2} U}{\partial z^{2}}$

% Laplacien d'un champ vectoriel : $\Delta \vec{A}=\left\lvert\, \begin{aligned} & \Delta A_{x}=\frac{\partial^{2} A_{x}}{\partial x^{2}}+\frac{\partial^{2} A_{x}}{\partial y^{2}}+\frac{\partial^{2} A_{x}}{\partial z^{2}} \\ & \Delta A_{y}=\frac{\partial^{2} A_{y}}{\partial x^{2}}+\frac{\partial^{2} A_{y}}{\partial y^{2}}+\frac{\partial^{2} A_{y}}{\partial z^{2}} \\ & \Delta A_{z}=\frac{\partial^{2} A_{z}}{\partial x^{2}}+\frac{\partial^{2} A_{z}}{\partial y^{2}}+\frac{\partial^{2} A_{z}}{\partial z^{2}}\end{aligned}\right.$

% Relations concernant les opérateurs mathématiques:

% $$
% \begin{aligned}
% & \overrightarrow{\operatorname{rot}}(\overrightarrow{\operatorname{rot}} \vec{Z})=\overrightarrow{\operatorname{grad}}(\operatorname{div} \vec{Z})-\Delta \vec{Z} \\
% & \Delta U=\operatorname{div}(\overrightarrow{\operatorname{grad} U)}
% \end{aligned}
% $$

% Trigonométrie :

% $$
% \begin{aligned}
% & \cos (a+b)=\cos a \cos b-\sin a \sin b \\
% & \cos (a-b)=\cos a \cos b+\sin a \sin b \\
% & \sin (a+b)=\sin a \cos b+\sin b \cos a \\
% & \sin (a-b)=\sin a \cos b-\sin b \cos a
% \end{aligned}
% $$

% $\cos (p)+\cos (q)=2 \cos \left(\frac{p+q}{2}\right) \cos \left(\frac{p-q}{2}\right)$

% $\cos (p)-\cos (q)=-2 \sin \left(\frac{p+q}{2}\right) \sin \left(\frac{p-q}{2}\right)$

% $\sin (p)+\sin (q)=2 \sin \left(\frac{p+q}{2}\right) \cos \left(\frac{p-q}{2}\right)$

% $\sin (p)-\sin (q)=2 \sin \left(\frac{p-q}{2}\right) \cos \left(\frac{p+q}{2}\right)$

% $\operatorname{ch}^{2}(x)-\operatorname{sh}^{2}(x)=1$
\Enonce{
On réalise, dans l'air, l'expérience des trous d'Young à l'aide du dispositif décrit et schématisé ci-dessous.
Un laser, de longueur d'onde dans le vide $\lambda$, émet un faisceau lumineux cylindrique d'axe $z^{\prime} z$.
On suppose par la suite, sauf mention contraire dans la question 1), que le faisceau du laser éclaire entièrement et de manière uniforme les différentes ouvertures qui sont placées sur son passage.
Une plaque opaque (P), percée de deux trous circulaires $S_{1}$ et $S_{2}$ de même taille et de faibles dimensions, est placée perpendiculairement à l'axe $z$'$z$.
On note $O^{\prime}$ le milieu du segment $\left[S_{1} S_{2}\right]$. Le point $O^{\prime}$ appartient à l'axe $z^{\prime} z$.
La distance entre les centres des deux trous $S_{1}$ et $S_{2}$ est notée $a$.
Le phénomène d'interférences est observé sur un écran (E) placé perpendiculairement à l'axe $z$ ' $z$. Soit $O$ le point de l'écran (E) appartenant à l'axe $z$'$z$.
La distance entre la plaque (P) et l'écran (E) est égale à $D$. On a ainsi $D=O^{\prime} O$.
L'espace est rapporté au repère cartésien $\left(O, \vec{e}_{x}, \vec{e}_{y}, \vec{e}_{z}\right)$ défini comme suit:

$\vec{e}_{z}$ : vecteur unitaire de l'axe $O z$, orienté de la plaque (P) vers l'écran (E).

$\vec{e}_{x}$ : vecteur unitaire de l'axe $O x$, parallèle à $\left[S_{1} S_{2}\right]$ et orienté de $S_{2}$ vers $S_{1}$.

$\vec{e}_{y}$ : vecteur unitaire de l'axe $O y$ tel que la base $\left(\vec{e}_{x}, \vec{e}_{y}, \vec{e}_{z}\right)$ soit orthonormée directe.

Dans tout le problème, l'indice de réfraction de l'air sera pris égal à 1.

%%%%%%%%%%%%%%%%%%%%%%%%%%%%%%%%%%%%%%%%%%%
\begin{fig}[label]
\centering
\includegraphics[width=0.7\columnwidth]{2024_02_01_d19b22528fcac7c6bd43g-2}
% \capt{}
\end{fig}
%%%%%%%%%%%%%%%%%%%%%%%%%%%%%%%%%%%%%%%
}

\subsection{Première partie : réalisation pratique du dispositif}

\Enonce{
\parag{Elargisseur de faisceau}

Nous nous plaçons, uniquement pour cette question, dans le cas où le laser dont nous disposons ne permet pas d'éclairer simultanément les deux trous $S_{1}$ et $S_{2}$ de la plaque (P): le diamètre du faisceau laser est inférieur à la distance $a$.

Nous souhaitons réaliser un dispositif permettant d'élargir le faisceau laser. Pour cela nous disposons de deux lentilles : une lentille divergente $L_{1}$ et une lentille convergente $L_{2}$ de distances focales respectives $f^{\prime}{ }_{1}$ et $f^{\prime}{ }_{2}$, telles que $\left|f^{\prime}{ }_{1}\right|<f^{\prime}{ }_{2}$.

On dispose les deux lentilles de telle sorte que le foyer image de la lentille divergente $L_{1}$ soit confondu avec le foyer objet de la lentille convergente $L_{2}$.

La lentille $L_{1}$ est placée avant la lentille $L_{2}$ sur le trajet du faisceau laser.

Soit $d_{a v}$ le diamètre du faisceau laser avant son élargissement et $d_{a p}$ le diamètre du faisceau laser après son élargissement.
}

\begin{quest}
  \QR{	 
  Faire un schéma du montage ainsi réalisé et mettant en évidence l'élargissement du faisceau incident.
  Calculer le rapport $\frac{d_{a p}}{d_{a v}}$ en fonction de $f^{\prime}{ }_{1}$ et $f^{\prime}{ }_{2}$.
  }{
    Le rapport entre les diamètres est donné par la relation du double afocal de lentilles :
\eqbox{
  \frac{d_{a p}}{d_{a v}}=-\frac{f_{2}^{\prime}}{f^{\prime}{ }_{1}}.
}
  }
  
  \QR{ Quelle est la couleur émise par le laser dans le cas d'un laser He-Ne de longueur d'onde dans le vide $\lambda=632,8 \mathrm{~nm}$ ?
  Quelles sont les longueurs d'onde et les couleurs respectives qui correspondent aux limites du spectre de la lumière visible?
  }{
    Il s'agit de la couleur \boxed{rouge} (laser usuel de TP)
  
  Le spectre visible correspond à l'intervalle de $400 \mathrm{~nm}$ (violet) à $750 \mathrm{~nm}$ (rouge).
  }
\end{quest}

\subsection{Deuxième partie : description qualitative du phénomène}

\Enonce{
\parag{Existence d'un champ d'interférences}
}

\begin{quest}
  \QR{	
   En se référant uniquement aux lois de l'optique géométrique, quelle devrait être l'allure de la figure observée sur l'écran (E)?
  
  Pour quelle raison l'optique géométrique ne permet-elle pas de prévoir l'existence d'un champ d'interférences dans le cas du dispositif des trous d'Young ?
  
  À quel phénomène physique doit-on faire appel pour en comprendre l'existence ?
  }{
    Dans le cadre de l'optique géométrique, on devrait voir deux petites taches lumineuses sur l'écran, images géométriques des trous.

    À cause du faible diamètre des trous, il y a un phénomène de \fbox{diffraction} de la lumière.
    
  }
  \QR{  Réaliser un schéma représentant le champ d'interférences.
  }{
    Voir cours.
  }
\end{quest}

\Enonce{
\parag{Description de la figure d'interférences}

Les réponses aux questions suivantes seront justifiées. Une démonstration quantitative ne sera toutefois pas exigée.
}

\begin{quest}
  \QR{	
  Qu'observe-t-on sur l'écran (E) ? Décrire précisément la figure d'interférences obtenue.
  }{
    Sur l'écran (E) on observe une alternance de taches sombres et brillantes selon l'axe Ox.  En effet, pour des raisons de symétrie, la différence de marche entre les deux ondes ne peut dépendre que de $x$, distance à l'écran suivant l'axe des trous.
    
  }
  
  \QR{  Qu'observe-t-on sur l'écran (E) si l'on obture l'un des deux trous?
  }{
    Si on obture le trou (S2), on observe la tache de diffraction (alternance anneaux sombres et brillants) centrée en $x=a / 2, y=z=0$.
  
  }
  
  \QR{  Comment est modifiée la figure d'interférences si on translate la plaque (P) suivant l'axe $O x$ ? suivant l'axe $O y$ ?
  }{
   Si on translate $(P)$ suivant $O x$, la figure de diffraction est translatée selon Ox. Idem selon Oy.
  }
  \QR{  Comment est modifiée la figure d'interférences si on translate l'écran (E) suivant l'axe $z$'$z$ ?
  }{
    Si on translate $(E)$ selon $z'z$, l'interfrange est modifié.
  }
  
\end{quest}

%%%%%%%%%%%%%
\subsection{Troisième partie : description quantitative du phénomène}

\Enonce{
\parag{Différence de chemin optique}

Soit un point $M$ de l'écran (E), de coordonnées $(x, y, 0)$ dans le repère $\left(O, \vec{e}_{x}, \vec{e}_{y}, \vec{e}_{z}\right)$.
}

\begin{quest}
  \QR{	
  Exprimer les coordonnées des trous $S_{1}$ et $S_{2}$ dans le repère $\left(O, \vec{e}_{x}, \vec{e}_{y}, \vec{e}_{z}\right)$.
  
  Exprimer les distances $S_{1} M$ et $S_{2} M$, respectivement entre les trous $S_{1}$ et $S_{2}$ et le point $M$. On exprimera $S_{1} M$ et $S_{2} M$ en fonction de $a, D, x$ et $y$.
  
  En déduire l'expression de la différence de chemin optique $\delta(M)=S_{2} M-S_{1} M$ au point $M$ entre les rayons issus de $S_{1}$ et $S_{2}$. On exprimera $\delta(M)$ en fonction de $a, D, x$ et $y$. Le calcul sera mené sans aucune approximation.
  }{ Les points sont de coordonnées  respectives $M(x, y, 0)$, $S_1(a/2,0,-D)$ et $S_2(-a/2,0,-D)$.
On a donc 
$$S_{1} M=\sqrt{(x-a / 2)^{2}+y^{2}+D^{2}},$$ et $$S_{2} M=\sqrt{(x+a / 2)^{2}+y^{2}+D^{2}}.$$
  }
  \QR{  La distance $a$ entre les deux trous étant petite par rapport à la distance d'observation $D$, et le point $M$ étant proche du point $\mathrm{O}$, on peut considérer que $a, x, y$ sont très petits devant $D$.
  
  En faisant un développement limité au premier ordre de l'expression de $\delta(M)$ obtenue précédemment, en déduire l'expression simplifiée de $\delta(M)$ en fonction de $a, D$ et $x$.
  }{
    Voir calcul du cours.
    \eqbox{
     \delta(M)=\frac{a  x}{D}.
    }
  }
  
  \QR{  En prenant en compte l'expression de $\delta(M)$ calculée à la question précédente, expliquer comment serait modifiée la figure d'interférences si on remplaçait les deux trous par deux fentes très fines appartenant à la plaque $(\mathrm{P})$, parallèles à l'axe $O y$ et distantes de $a$ ?
  }{
    La différence de marche ne dépend que de $x$; les franges d'interférences sont des lieux d'équation $x=\cte$. Si on remplace les trous par des fentes, la figure d'interférences conserve la même forme mais sera \fbox{plus lumineuse}.
    }
\end{quest}
  
\Enonce{
  \parag{Intensité lumineuse de l'onde résultante}
  On représente par $s_{1}(t)=s_{2}(t)=s_{0} \cos \textstyle\left(\frac{2 \pi c}{\lambda} t\right)$ l'expression des ondes respectivement aux points $S_{1}$ et $S_{2}$, où $s_{0}$ représente l'amplitude de l'onde considérée, $c$ représente la célérité de la lumière dans le vide et $t$ le temps.
  On néglige l'atténuation de l'onde entre les trous et le point $M$.
}

  \begin{quest}
    \QR{	
     Déterminer l'expression $s_{1 M}(t)$ de l'onde issue du trou $S_{1}$ lorsqu'elle arrive au point $M$. On exprimera $s_{1 M}(t)$ en fonction de $s_{0}, S_{1} M, c, \lambda$ et $t$.
    
    Déterminer, de même, l'expression $s_{2 M}(t)$ de l'onde issue du trou $S_{2}$ lorsqu'elle arrive au point $M$. On exprimera $s_{2 M}(t)$ en fonction de $s_{0}, S_{2} M, c, \lambda$ et $t$.
    }{
      Il s'agit de l'expression d'une vibration lumineuse monochromatique. Le temps de parcours dans le vide entre $S$ et $M$ vaut $\tau_{SM} = \frac{SM}{c}$.
      Les ondes s'écrivent donc $s_{1M}(t) = s_1(t-\tau_{S_1 M})$ et $s_{2M}(t) = s_2(t-\tau_{S_2 M})$, soit
      \eqbox{
        s_{1 M}(t)=s_{0} \cos \left(\frac{2 \pi c}{\lambda}\left(t-\frac{S_{1} M}{c}\right)\right) \text{ et } s_{2 M}(t)=s_{o} \cos \left(\frac{2 \pi c}{\lambda}\left(t-\frac{S_{2} M}{c}\right)\right).
      }
    }
    \QR{  En déduire l'expression $s_{M}(t)$ de l'onde qui résulte de la superposition des deux ondes $s_{1 M}(t)$ et $s_{2 M}(t)$ au point $M$. On exprimera $s_{M}(t)$ en fonction de $s_{0}, S_{1} M, S_{2} M, c, \lambda$ et $t$.
    Mettre l'expression de $s_{M}(t)$ sous la forme du produit d'un terme indépendant du temps (amplitude de l'onde) et d'un terme dépendant du temps.
    }{\udl{Principe de superposition} :
      $s(M)=s_{1 M}(t)+s_{2 M}(t)$. 
      
      D'après la relation trigonométrique $\cos(a)+\cos(b)=2 \cos \left(\frac{a+b}{2}\right) \cos \left(\frac{a-b}{2}\right)$,
      \eqbox{
        s_{M}(t)=2 s_{o} \cos \left(\frac{\pi c}{\lambda}\left(2 t-\frac{S_{1} M}{c}-\frac{S_{2} M}{c}\right)\right)  \cos \left(\frac{\pi}{\lambda}\left(S_{2} M-S_{1} M\right)\right).
      }
    }
    
    \QR{  Sachant que l'intensité lumineuse $I_{M}$ (appelée aussi éclairement) qui résulte, au point $M$, de l'onde $s_{M}(t)$ est proportionnelle au carré de l'amplitude de $s_{M}(t)$ avec $K$ constante de proportionnalité, exprimer l'intensité lumineuse $I_{M}$ au point $M$ en fonction de $s_{0}, K, \delta$ et $\lambda$ puis en fonction de $s_{0}, K, a, x, \lambda$ et $D$.
    }{
      \udl{Définition} : l'intensité lumineuse vaut $I_M = K \moy{s_M^2(t)}$. 

      On a donc
      $$I_M=4 K s_0^{2} \cos ^{2}\left(\frac{\pi}{\lambda}\left(S_{2} M-S_{1} M\right)\right)=4 K s_{0}^{2} \cos ^{2}\left(\frac{\pi}{\lambda} \delta\right),$$
      soit 
      \eqbox{
       I_M = 4 K s_{0}^{2} \cos ^{2}\left(\frac{\pi a x}{\lambda D}\right).
      }
    }
    
    \QR{  Calculer, en détaillant clairement le raisonnement effectué, l'expression de l'interfrange $i$ de la figure d'interférences. Exprimer $i$ en fonction de $a, \lambda$ et $D$.
    }{
      L'interfrange $i$ est la période (spatiale) de la fonction $x \mapsto I_M$, soit 
      \eqbox{
        i=\frac{\lambda {D}}{{a}}.
      }
    }
    
    \QR{  Tracer l'allure du graphe de $I_{M}$ en fonction de $x$.
    }{
      Voir cours.
    }
        \QR{  Quelle est la position de la frange d'ordre $0$ ?
    }{
      On définit l'ordre d'interférence par $\delta=p \lambda=\frac{a x_{p}}{D}$. L'ordre d'interférence $p=0$ correspond à $x=0$. La frange d'ordre $0$ est donc au centre de l'écran.
    }
  \end{quest}

\subsection{Quatrième partie : modification du dispositif interférentiel}
Nous nous plaçons encore pour cette partie dans le cas où la distance $a$ entre les deux trous étant petite par rapport à la distance d'observation $D$, et le point $M$ étant proche du point $\mathrm{O}$, on peut considérer que $a, x, y$ sont très petits devant $D$.

\Enonce{
\parag{Interposition d'une lame à faces parallèles}
Dans cette question uniquement, on rajoute devant le trou $S_{2}$ une petite lame (L) (verre ou mica) à faces parallèles, d'épaisseur $e$ et d'indice $n$ pour la longueur d'onde $\lambda$ utilisée. Le faisceau laser arrive toujours perpendiculairement à la plaque (P) et traverse la lame (L) sous incidence normale.

%%%%%%%%%%%%%%%%%%%%%%%%%%%%%%%%%%%%%%%%%%%
\begin{fig}[label]
\centering
\includegraphics[width=0.7\columnwidth]{2024_02_01_d19b22528fcac7c6bd43g-5}
% \capt{}
\end{fig}
%%%%%%%%%%%%%%%%%%%%%%%%%%%%%%%%%%%%%%%
}


\begin{quest}
  \QR{	
  Calculer la différence de chemin optique $\delta^{\prime}(M)$ au point $M$ entre les rayons issus de $S_{1}$ et $S_{2}$. Exprimer $\delta^{\prime}(M)$ en fonction de $n, e, a, x$ et $D$.
  
  En déduire, en utilisant les résultats de la partie précédente, de quelle distance $d$ la figure d'interférences sera translatée. On exprimera $d$ en fonction de $n, e, a$ et $D$.
  
  Dans quel sens se déplace la figure d'interférences par rapport au cas de la question précédente?
  }{
    La lame placée devant $S_{2}$ introduit une différence de marche $(n-1)e$ entre $S_{1}$ et $S_{2}$. La nouvelle différence de marche vaut $\delta^{\prime}(M)=\frac{a x}{D}+(n-1) e$
    
    L'interfrange n'est pas modifié puisque la différence de marche supplémentaire est indépendante de $x$. 
    
    La frange d'ordre 0 est telle que $\delta^{\prime}(M)=0$ donc 
    \eqbox{
     x_{0}=-\frac{(n-1) e D}{a}.
    }
    La figure est translatée de $d$. On a $n>1$ donc $d<0$ : les franges sont déplacées \guill{vers le bas}.
  }
  \end{quest}
  
\parag{Inclinaison du faisceau laser}

    Les rayons du faisceau laser ne sont plus parallèles à l'axe $z^{\prime} z$. Ils sont inclinés d'un angle $\alpha$ par rapport à cet axe. On se placera dans le cas où l'angle $\alpha$ est petit.
    
    % %%%%%%%%%%%%%%%%%%%%%%%%%%%%%%%%%%%%%%%%%%%
    % \begin{fig}[label]
    %   \centering
    %   \includegraphics[width=0.7\columnwidth]{2024_02_01_d19b22528fcac7c6bd43g-6}
    %   \capt{}
    % \end{fig}
    % %%%%%%%%%%%%%%%%%%%%%%%%%%%%%%%%%%%%%%%
    
    
    %%%%%%%%%%%%%%%%%%%%%%%%%%%%%%%%%%%%%%%%%%%
    \begin{fig}[label]
      \centering
      \includegraphics[width=0.7\columnwidth]{2024_02_01_d19b22528fcac7c6bd43g-6(1)}
      % \capt{}
    \end{fig}
    %%%%%%%%%%%%%%%%%%%%%%%%%%%%%%%%%%%%%%%
    \begin{quest}
      \QR{	
    Comment est modifiée la figure d'interférences ?
    
    Quelle est dans ce cas la position de la frange d'ordre $0$ ?
    }{
      On a de nouveau un déphasage entre $S_{1}$ et $S_{2}$.
      
      
      \udl{Théorème de Malus} : on construit le point $H$, projeté orthogonal de $S_2$ sur le rayon passant par $S_1$, qui est donc de même phase. Par la formule du déphasage, la différence de marche introduite vaut $HS_1$.


      \udl{Trigonométrie} : on a $HS_1 = S_2H \sin(\alpha) = a \sin(\alpha)$.
      La nouvelle différence de marche vaut donc $\delta^{\prime \prime}(M)=\frac{a x}{D}-a \sin (\alpha)$. 
      
      La figure d'interférences est donc translatée vers le haut d'une distance
      \eqbox{
      d=D \sin (\alpha).
      }
      La frange d'ordre 0 est telle que $\delta^{\prime \prime}(M)=0$, donc $x=d$.
    }
  \end{quest}
  
  
  
\end{exo}
%%%%%%%%%%%%%%%%%%%%%%%%%%%%%%%%%%%
% \subsection{Durée: $\mathbf{4}$ heures}
% N.B. : le candidat attachera la plus grande importance à la clarté, à la précision et à la concision de la rédaction. Si un candidat est amené à repérer ce qui peut lui sembler être une erreur d'énoncé, il le signalera sur sa copie et devra poursuivre sa composition en expliquant les raisons des initiatives qu'il a été amené à prendre.

% \subsection{RAPPEL DES CONSIGNES}
% \begin{itemize}
%   \item Utiliser uniquement un stylo noir ou bleu foncé non effaçable pour la rédaction de votre composition ; d'autres couleurs, excepté le vert, bleu clair ou turquoise, peuvent être utilisées, mais exclusivement pour les schémas et la mise en évidence des résultats.
%   \item Ne pas utiliser de correcteur.
%   \item Écrire le mot FIN à la fin de votre composition.
% \end{itemize}

%%%%%%%%%%%%%%%%%%%%%%%%%%%%%%%%%%%
\begin{exo}[PC][1][devoir]{Interférences et diffraction}
\Annale[Physique MP]{CCINP}{2024}
\Enonce{
Ce sujet s'intéresse à différents aspects des phénomènes ondulatoires comme les interférences et la diffraction ainsi qu'aux aspects corpusculaires associés. Dans l'étude des propriétés ondulatoires de diffraction, nous évoquerons l'effet Talbot.

Découvert en 1836 par Henry James Talbot, cet effet reste relativement peu connu. Ce photographe déposa un brevet en 1842 de \guill{multiplication d'une image} sans utilisation d'un système optique (auto-image). L'effet est redevenu d'actualité par ses applications dans les domaines artistique, médical et celui de la détection infrarouge.

Quand une onde plane incidente est envoyée sur un réseau diffractant plan placé dans le plan $z=0$, on peut observer des images de celui-ci (auto-images) régulièrement espacées. La distance entre deux auto-images successives est appelée longueur de Talbot $z_{T}$. Des images \guill{décalées} du réseau sont aussi observées à des distances régulières $z_{T} / 2, z_{T} / 4, z_{T} / n \ldots$ de l'auto-image formant ce qu'on nomme un tapis de Talbot (\Fig{talbot}).

\bonus{
Le sujet est constitué de quatre parties indépendantes, excepté la question Q10.f de la partie III qui est liée à la question Q9.e de la partie II.

Les parties I et IV font appel à la physique quantique.\\
La partie II fait appel aux interférences optiques à 2 et 3 ondes.\\
La partie III étudie certains aspects des réseaux, en optique et en électromagnétisme.
}

% Les données sont placées en début de sujet.\\
%%%%%%%%%%%%%%%%%%%%%%%%%%%%%%%%%%%%%%%%%%%
\begin{fig}[talbot]
\centering
\includegraphics[width=0.5\columnwidth]{2024_11_25_be97bbde3376be1663e6g-03}
\capt{Tapis de Talbot}
\end{fig}
%%%%%%%%%%%%%%%%%%%%%%%%%%%%%%%%%%%%%%%

% \href{https://upload.wikimedia.org/wikipedia/commons/thumb/6/69/Optical}{https://upload.wikimedia.org/wikipedia/commons/thumb/6/69/Optical} Talbot Carpet.png/460pxOptical Talbot Carpet.png

\parag{Données}
\begin{itemize}
  \item Constante de Planck $h=6,63 \cdot 10^{-34} \mathrm{~J} \cdot \mathrm{~s}$
  \item Constante de Planck réduite $\hbar=\frac{h}{2 \pi}=1,05 \cdot 10^{-34} \mathrm{~J} \cdot \mathrm{~s}$
  \item Vitesse de la lumière dans le vide $c=3,00 \cdot 10^{8} \mathrm{~m} \cdot \mathrm{~s}^{-1}$
  \item Constante de Boltzmann $k_{B}=1,38 \cdot 10^{-23} \mathrm{~J} \cdot \mathrm{~K}^{-1}$
  \item Masse de l'électron $m=9,11 \cdot 10^{-31} \mathrm{~kg}$
  \item Charge élémentaire $e=1,6 \cdot 10^{-9} \mathrm{C}$
  \item Constante des gaz parfaits $R=8,31 \mathrm{~J} \cdot \mathrm{~K}^{-1} \cdot \mathrm{~mol}^{-1}$
  \item Constante d'Avogadro $N_{A}=6,02 \cdot 10^{23} \cdot \mathrm{~mol}^{-1}$
\end{itemize}

\parag{Notations}
\begin{itemize}
  \item Le nombre imaginaire pur racine carrée de -1 est noté $j$.
  \item La notation $\langle f(M, t)\rangle$ représente la moyenne temporelle de la fonction $f(M, t)$.
  \item La fonction sinus cardinal est par définition $\text{sinc}(u)=(\sin u) / u$.
  \item Elle possède un maximum principal égal à 1 pour $u=0$, des zéros pour les valeurs de $u$ égales à $u_{\text {max }}=n \pi$ avec $n$ entier relatif non nul et des maximas secondaires pour des valeurs proches de $u_{\text {min }}=(2 p+1) \pi / 2$ avec $p$ entier relatif différent de 0 et -1.
  \item Les deux premiers maxima secondaires de la fonction (sinus cardinal) valent : 0,047 et 0,016.
  \item Interfrange dans un système à division du front d'onde de deux fentes d'Young : on pourra utiliser, sans démonstration, que celui-ci vaut
  \begin{equation}
    i =\frac{\lambda D}{a},
  \end{equation}
  où $D$ est la distance entre le plan des fentes et l'écran, $a$ la distance entre les fentes et $\lambda$ la longueur d'onde de la source. Les rayons sont considérés paraxiaux (faiblement inclinés par rapport à l'axe du système optique).
\end{itemize}
}


\bonus{
\subsubsection{Série de Fourier d'une fonction \guill{créneaux}}
Pour la fonction périodique \guill{carré}de hauteur 1, de largeur $\varepsilon$, de période $a$ et de valeur moyenne $\textstyle\frac{1}{2}$, on peut écrire $f(x)=\sum_{n=-\infty}^{n=+\infty} C_{n} \exp \left(-j n \frac{2 \pi x}{a}\right)$ avec $C_{0}=\frac{\pi x}{a}$ et $C_{n}=\frac{\sin \left(\frac{n \pi x}{a}\right)}{n \pi}$.\\
%%%%%%%%%%%%%%%%%%%%%%%%%%%%%%%%%%%%%%%%%%%
\begin{fig}[label]
\centering
\includegraphics[width=0.5\columnwidth]{2024_11_25_be97bbde3376be1663e6g-04}
\capt{Fonction créneau.}
\end{fig}
%%%%%%%%%%%%%%%%%%%%%%%%%%%%%%%%%%%%%%%
}

\bonus{
\subsection{Dualité onde-corpuscule pour la lumière}
\begin{quest}
  \item a) On considère un faisceau parallèle de photons associés à une onde électromagnétique, de longueur d'onde $\lambda$, se propageant dans l'air assimilé au vide dans la direction Oz de vecteur unitaire $\overrightarrow{e_{z}}$ (figure 1).\\
  Rappeler ce que vaut leur énergie $E$ en fonction de $\lambda, c$ et de $h$ et leur quantité de mouvement $\vec{p}$ en fonction de $\lambda, h$ et de $\overrightarrow{e_{z}}$.\\
  b) Calculer l'énergie en eV d'un photon de lumière bleue de longueur d'onde 475 nm .
  
  \item Ce faisceau parallèle cylindrique de rayon $R$ arrive face à un écran, perpendiculaire à l'axe du faisceau, percé d'un trou circulaire $T_{1}$ de centre $O$ et de rayon $r$ (inférieur à $R$ ) (figure 1).\\
\end{quest}
  %%%%%%%%%%%%%%%%%%%%%%%%%%%%%%%%%%%%%%%%%%%
\begin{fig}[label]
\centering
\includegraphics[width=0.5\columnwidth]{2024_11_25_be97bbde3376be1663e6g-05}
\capt{}
\end{fig}
%%%%%%%%%%%%%%%%%%%%%%%%%%%%%%%%%%%%%%%

Figure 1 - Géométrie du dispositif à un trou\\
a) Établir, à partir de l'inégalité d'Heisenberg spatiale, qu'il y a forcément ouverture angulaire du faisceau.\\
b) Donner un ordre de grandeur de cette ouverture angulaire supposée petite. Commenter. Comment appelle-t-on ce phénomène ?

\begin{quest}
  \item Citer une expérience qui met en évidence l'aspect corpusculaire de la lumière.\\
  \item Un électron est expulsé d'un métal sous l'effet d'une radiation s'il absorbe une énergie au moins égale à We, énergie appelée travail " d'extraction ".
\end{quest}
  
Le tableau ci-dessous indique les valeurs du travail " d'extraction " pour différents métaux :

\begin{tabular}{|c|c|c|c|c|c|}
\hline
Métal & Cs & Na & K & Ti & Fe \\
\hline
We en eV & 1,15 & 2,11 & 2,22 & 4,33 & 4,67 \\
\hline
\end{tabular}

Avec quels métaux cités dans le tableau, la lumière bleue du Q1.b permet-elle d'obtenir un effet photoélectrique ? Justifier.\\
Quelle sera la vitesse maximale des électrons émis ? Commenter.\\
\begin{quest}
  \item a) On utilise une source optique de puissance 1 mW : évaluer l'ordre de grandeur du nombre de photons qui sortent de la source par unité de temps en supposant le faisceau rigoureusement monochromatique de longueur d'onde $\lambda=632,8 \mathrm{~nm}$.\\
  b) À quel niveau de puissance faudrait-il descendre, pour une source, pour que la lumière qu'elle émet soit détectée photon par photon? On admet que les détecteurs photoniques ont un temps de réponse de l'ordre de la picoseconde (le préfixe pico correspond à $10^{-12}$ ).\\
  c) Décrire l'évolution des observations sur le détecteur en fonction de la durée d'observation dans la situation des interférences, par deux fentes d'Young éclairées par la source du Q5.b, modélisée par une source à photons uniques.\\
  d) Traduire ces observations en termes d'amplitude de probabilité pour un photon, détecté en un point M. Le principe de superposition s'applique-t-il aux probabilités ou aux amplitudes de probabilités?
\end{quest}
}

\bonus{
\subsection{Interférences lumineuses à $2$ ou $3$ ondes}
}

\Enonce{
On considère (\Fig{deux_trous}) un faisceau de lumière parallèle de longueur d'onde $\lambda$, se propageant dans la direction $(Oz)$. Ce faisceau arrive sur un écran placé dans le plan $(xOy)(z=0)$ percé de deux trous identiques $T_{1}$ et $T_{2}$. Les centres des trous $\rm O_{1}$ et $\rm O_{2}$ ont pour coordonnées respectivement $\textstyle\left(\frac{a}{2}, 0,0\right)$ et $\textstyle\left(-\frac{a}{2}, 0,0\right)$. Le rayon des trous est de l'ordre de grandeur de la longueur d'onde. Ceci permet de supposer qu'il existe un champ d'interférences qui est la zone commune aux deux faisceaux diffractés par les trous. On modélise chaque trou par une source secondaire ponctuelle émettant une lumière uniforme dans le champ d'interférences. Ces sources secondaires sont cohérentes entre elles.\\
%%%%%%%%%%%%%%%%%%%%%%%%%%%%%%%%%%%%%%%%%%%
\begin{fig}[deux_trous]
\centering
\includegraphics[width=0.5\columnwidth]{2024_11_25_be97bbde3376be1663e6g-06}
\capt{Géométrie du dispositif à deux trous.}
\end{fig}
%%%%%%%%%%%%%%%%%%%%%%%%%%%%%%%%%%%%%%%
}

\begin{quest}
  \item On observe sur un écran placé dans le plan $z=D$, en un point $M$ de coordonnées $(x, y, D)$. On suppose que $D$ est très grand devant $a,|x|$ et $|y|$. Le montage est réalisé dans l'air d'indice égal à l'unité.
  \begin{enumerate}
    \QR{ De quel type de division interférentielle s'agit-il ? Les interférences sont-elles localisées ?
    }{
      Le faisceau lumineux est séparé en deux par les deux fentes, on a donc une division interférentielle du \fbox{front d'onde}.  
      
      
      Les interférences s'observent dans une région étendue de l'espace (le champ d'interférences), elles sont donc \fbox{non localisées}.
    }
    \QR{ Établir (dans le cadre de l'approximation scalaire de l'optique) l'expression de la différence de marche $\delta(M)$. Établir l'expression de l'intensité $I(M)$ au point $M$ en notant $I_{\rm max}$ l'intensité maximale.
    }{
      Avec l'hypothèse \(D \gg x\) et \(D \gg a\), on peut considérer les deux rayons parallèles entre eux.

      On trace alors la surface d'onde partant de \(O_{1}\) et coupant le rayon \(O_{2} \mathrm{M}\) au point H.
      
      La différence de marche est alors la distance \(\mathrm{O}_{2} \mathrm{H}\).
      Par trigonométrie on trouve : \(\delta=a \sin \theta\).
      En considérant le rayon moyen, par trigonométrie, on a : 
      % \eqbox{
      $$\tan\theta = \frac{x}{D}.$$
      % }
      Le faisceau lumineux est séparé en deux par les deux fentes, on a donc une
\begin{center}
  \includegraphics[width=0.7\columnwidth]{2024_12_04_e0c7d5f4172e8f1d9be1g-03(1)}
\end{center}      
      En faisant l'approximation des petits angles, on trouve alors : \[\delta=\frac{a x}{D}.\]
      Les deux sources secondaires \(O_{1}\) et \(O_{2}\) proviennent de la même source principale par division du front d'onde, elles sont donc cohérentes et de même amplitude.
      On peut alors appliquer la formule de Fresnel : $I(M)=2 I_{0}\left(1+\cos \left(\frac{2 \pi \delta}{\lambda}\right)\right)$, soit 
      \eqbox{
       I(M) = \frac{I_{\max }}{2}\left[1+\cos \left(\frac{2 \pi a x}{\lambda D}\right)\right].
      }
      Ou encore : \(I(M)=2 I_{0}(1+\cos (\varphi))\) avec \(\varphi=2 \pi \frac{a x}{\lambda D}\).
    }
  \end{enumerate}
    
  \item 
  \begin{enumerate}
    \QR{ 
    Décrire ce qu'on doit voir sur l'écran dans le cadre de ces hypothèses et exprimer l'interfrange $i$ en fonction de la fréquence d'émission $\nu$, de $c$ et des paramètres géométriques du dispositif.

    }{
      Sur l'écran on observe la figure d'interférence de trous d'Young, c'est-à-dire des franges verticales régulièrement espacées dans la tache de diffraction, représentée ci-contre.\\
      L'interfrange \(i\) est la périodicité spatiale de l'intensité calculée dans la question précédente, c'est donc la périodicité du terme \(\cos \left(\frac{2 \pi a x}{\lambda D}\right)\), on a donc: 
      \eqbox{
      i=\frac{\lambda D}{a}=\frac{c D}{a \nu}.
      }
      \begin{center}
        \includegraphics[width=0.7\columnwidth]{2024_12_04_e0c7d5f4172e8f1d9be1g-03}
      \end{center}
      
    }
    \QR{ Pourquoi, dans ce cadre, peut-on remplacer les deux trous par deux fentes fines identiques parallèles à $(Oy)$ ? Quel en est l'intérêt? La figure est-elle transformée si on translate de façon \guill{raisonnable} en bloc les fentes dans leur plan?
    }{
      Les fentes d'Young se comportent comme un ensemble de trous d'Young accolés les uns aux autres. Les différences de marches changent pour chaque couple de trous, mais entrainent la même figure d'interférences.

      Ainsi, les figures se superposant, on retrouve la même figure d'interférence mais bien \fbox{plus lumineuse}.

      La différence entre les fentes d'Young et les trous d'Young provient uniquement de la figure de diffraction qui a une symétrie rectiligne au lieu de circulaire.
    }
  \end{enumerate}
    
\end{quest}

\Enonce{
La source est en réalité quasi-monochromatique à profil spectral \guill{rectangulaire} de largeur $\Delta \lambda$ autour de $\lambda_{0}$, avec $\Delta \lambda \ll \lambda_{0}$. Ce profil spectral, en fonction de la fréquence d'émission, est représenté sur la \Fig{profil_spectre}. On admet que l'intensité émise par une bande spectrale de largeur $d \nu$ autour de $\nu$ vaut $d I=J_{\nu} d \nu$.

  %%%%%%%%%%%%%%%%%%%%%%%%%%%%%%%%%%%%%%%%%%%
  \begin{fig}[profil_spectre]
    \centering
\includegraphics[width=0.7\columnwidth]{2024_11_25_be97bbde3376be1663e6g-07}
\capt{Profil \guill{rectangulaire} d'une source quasi-monochromatique. Par commodité de représentation, l'échelle n'est pas respectée $\left(\Delta \nu \ll \nu_{0}\right)$.}
\end{fig}
%%%%%%%%%%%%%%%%%%%%%%%%%%%%%%%%%%%%%%%
}

\begin{quest}
  \item
  \begin{enumerate}
    \QR{ Établir l'expression de l'intensité $I(M)$ au point $M$ en notant $I_{\rm max}$ l'intensité maximale et montrer qu'elle peut s'écrire sous la forme 
    \begin{equation}
      I(M)=\frac{I_{\max }}{2}\left(1+V(M) \cos \left( \frac{2 \pi x}{i} \right)\right),
    \end{equation}
     avec $i$ l'interfrange correspondant à la valeur centrale de la raie. On rappelle que $\sin p-\sin q=2 \sin \left( \frac{p-q}{2} \right) \cos \left( \frac{p+q}{2} \right)$.

     }{
      L'intensité élémentaire s'écrit : \(d I=J_{\nu} d \nu=\frac{I_{0}}{\Delta \nu} d \nu\).
      On reprend le calcul d'intensité précédent : \[d I(M)=\frac{I_{0}}{\Delta \nu} d \nu\left(1+\cos \left(\frac{2 \pi \delta}{\lambda}\right)\right)=\frac{I_{0}}{\Delta \nu} d \nu\left(1+\cos \left(\frac{2 \pi \delta}{c} \nu\right)\right)\].
      L'intensité totale s'écrit donc: \[I(M)=\int_{\nu_{0}-\frac{\Delta \nu}{2}}^{\nu_{0}+\frac{\Delta \nu}{2}} d I(M)=\frac{I_{0}}{\Delta \nu} \int_{\nu_{0}-\frac{\Delta \nu}{2}}^{\nu_{0}+\frac{\Delta \nu}{2}}\left(1+\cos \left(\frac{2 \pi \delta}{c} \nu\right)\right) d \nu\].
      D'où : \[I(M)=\frac{I_{0}}{\Delta \nu}\left[\Delta \nu+\frac{c}{2 \pi \delta}\left[\sin \left(\frac{2 \pi \delta}{c}\left(\nu_{0}+\frac{\Delta \nu}{2}\right)\right)-\sin \left(\frac{2 \pi \delta}{c}\left(\nu_{0}-\frac{\Delta \nu}{2}\right)\right)\right]\right]\].
      En utilisant la formule de trigonométrie fournie : \[I(M)=\frac{I_{0}}{\Delta \nu}\left[\Delta \nu+\frac{c}{\pi \delta} \sin \left(\frac{\pi \delta}{c} \Delta \nu\right) \cos \left(\frac{2 \pi \delta}{c} \nu_{0}\right)\right]\].
      Et donc finalement : \(I(M)=I_{0}\left[1+\frac{c}{\pi \delta \Delta \nu} \sin \left(\frac{\pi \delta}{c} \Delta \nu\right) \cos \left(\frac{2 \pi x}{i}\right)\right]\) avec \(i=\frac{c}{\delta \nu_{0}}=\frac{c D}{a \nu_{0}}\).
      Et donc:
      \eqbox{
      I(M)=\frac{I_{\max }}{2}\left[1+\operatorname{sinc}\left(\frac{\pi \delta}{c} \Delta \nu\right) \cos \left(\frac{2 \pi x}{i}\right)\right].
      }
     }
    \QR{ Exprimer la visibilité $V(M)$ correspondante.
    }{
      En comparant le résultat à l'expression fournie, on obtient la visibilité (le contraste) : 
      \eqbox{
        V(M)=\operatorname{sinc}\left(\frac{\pi \delta}{c} \Delta v\right).
      }
    }
    \QR{ Exprimer en fonction de $\Delta \lambda$ et $\lambda_{0}$ la longueur de cohérence $L$, c'est-à-dire la plus petite valeur de la différence de marche $\delta$ à partir de laquelle les franges ne sont plus visibles. Vérifier que ce résultat correspond au critère de brouillage des franges portant sur l'ordre d'interférences.
    On rappelle le lien entre les deux largeurs spectrales $\Delta \nu=\frac{c}{\lambda_{0}^{2}} \Delta \lambda$.

    }{
      La longueur de cohérence \(L\) correspond à la différence de marche produisant la première annulation de la visibilité.
      On a donc la première annulation de \(\sin \left(\frac{\pi \delta}{c} \Delta v\right)\) pour une différence de marche non nulle : \(\frac{\pi L}{c} \Delta v=\pi\).
      On a donc: 
      \eqbox{
        L=\frac{c}{\Delta v}.
      } 
      On calcule l'ordre d'interférence correspondant : \(p=\frac{L}{\lambda}=\frac{c}{\Delta v \lambda}\).
      On calcule l'écart d'ordre d'interférence : \(\Delta p=p_{0}-p=\frac{c}{\Delta v}\left(\frac{1}{\lambda_{0}}-\frac{1}{\lambda_{0}+\frac{\Delta \lambda}{2}}\right)=\frac{c}{\Delta v} \frac{\frac{\Delta \lambda}{2}}{\lambda_{0}\left(\lambda_{0}+\frac{\Delta \lambda}{2}\right)} \simeq \frac{c}{\Delta v} \frac{\Delta \lambda}{2 \lambda_{0}^{2}}\), soit
      \eqbox{
       \Delta p \simeq \frac{1}{2}.
      } 
      On retrouve bien un écart d'ordre d'interférence de \(\frac{1}{2}\), ce qui correspond bien à la condition de brouillage.
    }
    \item Établir l'expression de la durée $\tau$ des trains d'onde ou temps de cohérence.
    \sol{
      La durée de cohérence est : 
      \eqbox{
        \tau=\frac{L}{c}=\frac{1}{\Delta v}.
      }
    }
    \QR{ Justifier pourquoi on définit le nombre d'interfranges visibles par $N=2 L / \lambda_{0}$.
    }{
      On ne peut observer d'interfrange que dans l'intervalle où l'on ne brouille pas.
      La différence de marche doit donc rester inférieure à la longueur de cohérence.
      On a donc: \(\delta_{\max }=\frac{a x_{\max }}{D}= \pm L\).
      D'où : 
      \eqbox{
      x_{\max }= \pm \frac{L D}{a}.
      }
      L'intervalle de positions sur lequel on peut observer les franges est donc: \(\Delta x=\frac{2 L}{a D}\).
      On trouve le nombre de franges visibles en utilisant l'interfrange : \(N=\frac{\Delta x}{i}=\frac{2 L D}{a} \frac{a}{\lambda_{0} D}\), soit 
      \eqbox{
      N=\frac{2 L}{\lambda_{0}}.
      }
      On retrouve bien l'expression indiquée dans le sujet.
    }
    \QR{ Dans le tableau ci-dessous sont indiquées des caractéristiques de sources quasi-monochromatiques. Après l'avoir recopié sur votre copie, le compléter et le commenter.
    }{
      On recopie et on complète le tableau fourni.

      \begin{center}
      \begin{tabular}{|c|c|c|c|c|c|}
      \hline
      Source & \(\lambda_{0}\) en nm & \(\Delta \lambda\) en nm & \(\tau\) en s & \(L\) en m & \(N\) \\
      \hline
      Laser He -Ne & 632,991 & 0,001 & \(1 \cdot 10^{-9}\) & \(4 \cdot 10^{-1}\) & \(1 \cdot 10^{6}\) \\
      \hline
      Raie rouge de l'hydrogène & 656,2 & 0,1 & \(1 \cdot 10^{-11}\) & \(4 \cdot 10^{-3}\) & \(1 \cdot 10^{4}\) \\
      \hline
      Lumière blanche filtrée & 500 & 20 & \(4 \cdot 10^{-14}\) & \(1 \cdot 10^{-5}\) & \(5 \cdot 10^{1}\) \\
      \hline
      \end{tabular}
      \end{center}
      
      On retrouve bien le résultat connu que plus une source est monochromatique plus elle est cohérente et plus on voit d'interfranges.
      
    }
  \end{enumerate}
  \end{quest}

  % \begin{table}[h!]
  \begin{center}
    \begin{tabular}{|l|c|c|c|c|c|}
      \hline
      Source & $\lambda_{0}$ en nm & $\Delta \lambda$ en nm & $\tau$ en s & $L$ en m & $N$ \\
      \hline
      Laser $\mathrm{He}-\mathrm{Ne}$ & 632,991 & 0,001 & $?$ & $?$ & $?$ \\
      \hline
      Raie rouge de l'hydrogène & 656,2 & 0,1 & $?$ & $?$ & $?$ \\
      \hline
      Lumière blanche filtrée & 500 & 20 & $?$ & $?$ & $?$ \\
      \hline
    \end{tabular}
    % \caption{C$aractéristiques de sources monochromatiques.}
  \end{center}
    % \label{}
  % \end{table}

\Enonce{
Dans la suite de cette partie, on considère que la source est rigoureusement monochromatique de longueur d'onde $\lambda$.  Une troisième fente, identique au deux autres, parallèle à $(Oy)$, est placée en $\mathrm{O}(0,0,0)$.
Le plan des trois fentes est placé orthogonalement à l'axe de révolution commun de deux lentilles minces convergentes de distance focale $f^{\prime}=50 \mathrm{~cm}$. La lumière provient d'une fente source monochromatique, parallèle aux fentes diffractantes, placée au foyer principal objet de la première lentille (lentille d'entrée). L'écran d'observation est confondu avec le plan focal image de la seconde lentille (lentille de sortie), placée en aval du plan percé.}

\begin{quest}
  \item
  \begin{enumerate}
    \QR{ 
     Faire une représentation schématique du montage.
    }{
      Voir question suivante.
     }
    \QR{Représenter les cheminements des trois rayons qui interfèrent en un point M de l'écran.
    }{
      On ajoute trois rayons qui interfèrent.\\
\begin{center}
  \includegraphics[width=0.7\columnwidth]{2024_12_04_e0c7d5f4172e8f1d9be1g-05}
\end{center}      
      % Remarque : nouveau problème du sujet, où la distance \(a\) défini toujours la distance entre les deux trous extrêmes alors qu'il aurait été plus cohérent (vu la suite du sujet) de garder a comme la distance entre deux fentes.
    }
    \QR{Établir la nouvelle fonction intensité $I(\varphi)$ en posant $\varphi=2 \pi \frac{a x}{\lambda f^{\prime}}$. On fera figurer l'intensité $I_{0}$ qui correspondrait à celle d'une fente unique identique.
    }{
      La différence de marche entre deux fentes successives se calcule comme précédemment : \[\delta=\frac{a x}{2 f^{\prime}}.\] Le déphasage entre deux fentes successives est donc : \[2 \pi \frac{\delta}{\lambda}=\pi \frac{a x}{\lambda f^{\prime}}=\frac{\varphi}{2}.\]
      Les trois vibrations sont donc:
      
      \[
      \begin{aligned}
        s_1(M, t)= & S_{0} \cos \left(\omega t+\frac{\varphi}{2}\right), \\
      & s_{2}(M, t)=S_{0} \cos (\omega t), \\
      & s_{3}(M, t)=S_{0} \cos \left(\omega t-\frac{\varphi}{2}\right) .
      \end{aligned}
      \]
      
      L'intensité est alors: \[I(\varphi)=K\left\langle\left(s_{1}+s_{2}+s_{3}\right)^{2}\right\rangle=K\left(\left\langle s_{1}^{2}\right\rangle+\left\langle s_{2}^{2}\right\rangle+\left\langle s_{3}^{2}\right\rangle+2\left\langle s_{1} s_{2}\right\rangle+2\left\langle s_{1} s_{3}\right\rangle+2\left\langle s_{2} s_{3}\right\rangle\right).\] Or : \[\left\langle s_{1}^{2}\right\rangle=\left\langle s_{2}^{2}\right\rangle=\left\langle s_{3}^{2}\right\rangle=\frac{s_{0}^{2}}{2}=I_{0}.\]

      On utilise la formule : \(\cos (a) \cos (b)=\frac{1}{2}[\cos (a+b)+\cos (a-b)]\).
      On a : \[\left\langle s_{1} s_{2}\right\rangle=\left\langle S_{0}^{2} \cos \left(\omega t+\frac{\varphi}{2}\right) \cos (\omega t)\right\rangle=\frac{S_{0}^{2}}{2}\left\langle\cos \left(2 \omega t+\frac{\varphi}{2}\right)\right\rangle+\frac{S_{0}^{2}}{2}\left\langle\cos \left(\frac{\varphi}{2}\right)\right\rangle=I_{0} \cos \left(\frac{\varphi}{2}\right).\]
      On a : \[\left\langle s_{1} s_{3}\right\rangle=\left\langle S_{0}^{2} \cos \left(\omega t+\frac{\varphi}{2}\right) \cos \left(\omega t-\frac{\varphi}{2}\right)\right\rangle=\frac{S_{0}^{2}}{2}\langle\cos (2 \omega t)\rangle+\frac{S_{0}^{2}}{2}\langle\cos (\varphi)\rangle=I_{0} \cos (\varphi).\]
      On a : \(\left\langle s_{2} s_{3}\right\rangle=\left\langle S_{0}^{2} \cos (\omega t) \cos \left(\omega t-\frac{\varphi}{2}\right)\right\rangle=\frac{S_{0}^{2}}{2}\left\langle\cos \left(2 \omega t-\frac{\varphi}{2}\right)\right\rangle+\frac{S_{0}^{2}}{2}\left\langle\cos \left(-\frac{\varphi}{2}\right)\right\rangle=I_{0} \cos \left(\frac{\varphi}{2}\right)\).
      On a donc au final: 
      \eqbox{
      I(\varphi)=3 I_{0}+4 I_{0} \cos \left(\frac{\varphi}{2}\right)+2 I_{0} \cos (\varphi).
      }
    }
    \QR{Représenter, sur un même graphe, les fonctions intensités $I(\varphi)$ pour les deux fentes  et pour les trois fentes de la question  en fonction de la variable $\varphi$ sur l'intervalle $[0,2 \pi]$.
    }{
      On réalise la représentation graphique demandée.

      Remarque : ici le tracé demandé pose soucis puisque la variable \(\varphi\) n'est pas la même dans les deux cas de figure. Il n'y a donc pas vraiment de comparaison possible.
\begin{center}
  \includegraphics[width=0.7\columnwidth]{2024_12_04_e0c7d5f4172e8f1d9be1g-06}
\end{center}    
}
    \QR{Commenter en comparant la luminosité et la largeur des zones les plus brillantes.
    }{
      On voit que la luminosité des zones brillantes est plus importante dans le cas à 3 fentes que dans le cas à 2 fentes.

      On note également que les zones brillantes sont un peu plus fine dans le cas à 3 fentes que dans le cas à 2 fentes.
    }
    \QR{On donne sur la photo 2 ce qu'on voit sur l'écran pour une lumière monochromatique rouge de longueur d'onde $\lambda=633 \mathrm{~nm}$. À partir de cette photo 2, évaluer la distance entre les deux fentes en considérant que seule la zone comprise entre les abscisses $x=2,1 \mathrm{~mm}$ et $x=-2,1 \mathrm{~mm}$ correspond au calcul de l'intensité tel qu'il a été fait, c'est-à-dire sans tenir compte de la largeur des fentes diffractantes.
    }{
      Sur l'interférogramme donné, on mesure que la largeur de la tache centrale est: \(2 x_{1}=1,0 \mathrm{~mm}\).
      On cherche la première annulation de l'intensité : \(I\left(\varphi_{1}\right)=0\).
      On calcule la dérivée de l'intensité : \(\frac{d I}{d \varphi}=-2 \sin \left(\frac{\varphi}{2}\right)-2 \sin (\varphi)\).
      La dérivée s'annule pour \(\frac{\varphi}{2}=-\varphi+2 n \pi\) avec \(n \in \mathbb{Z}\).
      D'où : \(\varphi=\frac{4 n \pi}{3}\).
      La première annulation pour \(\varphi \neq 0\) est donc : \(\varphi_{1}=\frac{4 \pi}{3}\).
      Or on a: \(\varphi_{1}=2 \pi \frac{a x_{1}}{\lambda f^{\prime}}\).
      D'où : 
      \eqbox{
      a=\frac{2 \lambda f^{\prime}}{3 x_{1}}.
      }
      \AN~:$a=\frac{2 \times 633 \cdot 10^{-9} \times 50 \cdot 10^{-2}}{3 \times 0,50 \cdot 10^{-3}}=4,4 \cdot 10^{-4} \mathrm{~m}=0,44 \mathrm{~mm}$.

      \udl{Remarque} : ici le \(\varphi_{1}=\frac{4 \pi}{3}\) trouvé dans cette question ne correspond pas au graphique précédent car le déphasage \(\varphi\) du graphique ne correspond pas au \(\varphi\) calculé ici.
    }
  \end{enumerate}
  \end{quest}
  %%%%%%%%%%%%%%%%%%%%%%%%%%%%%%%%%%%%%%%%%%%
  \Enonce{
  \begin{fig}[label]
\centering
\includegraphics[width=0.5\columnwidth]{2024_11_25_be97bbde3376be1663e6g-08}
\capt{Interférences à 3 fentes (\href{http://plateforme.sillages.info}{plateforme.sillages.info}). Les graduations sont en mm.}
\end{fig}
%%%%%%%%%%%%%%%%%%%%%%%%%%%%%%%%%%%%%%%
  }

\bonus{

\subsection{Diffraction par un réseau et effet Talbot en optique}
Q10. On se place dans le cas de $N$ fentes identiques parallèles à Oy (figure 4) équidistantes (période a). La source ponctuelle S , supposée rigoureusement monochromatique de longueur d'onde $\lambda=632,8 \mathrm{~nm}$ est placée au foyer objet d'une lentille mince convergente de distance focale $f^{\prime}$. On observe en sortie sur un écran placé dans le plan focal image d'une seconde lentille mince de distance focale $f^{\prime}$ et de centre optique O . On note Oz l'axe optique commun des deux lentilles.\\
a) Faire une représentation du système optique et représenter le cheminement d'un rayon depuis la source jusqu'en un point $M\left(x, y, f^{\prime}\right)$ de l'écran d'observation.\\
On pose $\varphi=2 \pi \frac{x a}{\lambda f^{\prime}}$.\\
b) Établir que la fonction intensité $I(\varphi)$ peut s'écrire sous la forme $I(\varphi)=N^{2} I_{0}\left(\frac{\sin N \frac{\varphi}{2}}{N \sin \frac{\varphi}{2}}\right)^{2}=N^{2} I_{0} f(\varphi)^{2}$ en notant $I_{0}$ l'intensité qui correspondrait à celle d'une fente unique identique.\\
c) En déduire les valeurs de $\varphi$ qui annulent le numérateur de la fonction $I(\varphi)$. En déduire les valeurs de $\varphi$ qui donnent des maxima d'intensité et celles qui donnent des zéros d'intensité. Combien existe-t-il de zéros d'intensité entre deux maximums d'intensité successifs?\\
d) Établir le rapport H de l'intensité des franges les plus brillantes à celle produite par une seule fente identique éclairée de la même façon.\\
e) Déterminer la « largeur» $\delta \varphi$ des franges brillantes.\\
f) Commenter en généralisant votre réponse à la question Q9.e.

Q11. Le résultat obtenu en Q10.c pour les maxima d'intensité, établi dans le cadre de l'approximation de Gauss, peut être réécrit sous la forme $\sin \theta_{n}=n \frac{\lambda}{a}$ avec $n$ entier relatif. Elle est appelée " formule du réseau " et elle reste valable pour des rayons non paraxiaux. Justifier, sans calcul de l'intensité, cette formule. Comment nomme-t-on l'entier $n$ ? Comment doit-on la modifier la formule si la situation n'est plus en incidence normale?

On suppose que l'onde électromagnétique incidente, associée à un rayon issu de la source ponctuelle à l'infini dans la direction $\overrightarrow{e_{z}}$, se propageant dans l'air d'indice unité est caractérisée par un champ électrique en notation complexe :

$$
\overrightarrow{E_{l}}=E_{0} \exp j\left(\omega t-\overrightarrow{k_{l}} \cdot \overrightarrow{O M}\right) \overrightarrow{e_{y}}
$$

avec $E_{0}$ et $\omega$ grandeurs constantes, uniformes et réelles et le vecteur d'onde $\vec{k}_{l}$ uniforme.\\
Q12. a) Indiquer les caractéristiques de cette onde incidente, en précisant sa polarisation, son caractère harmonique, sa planéité et l'expression du vecteur $\overrightarrow{k_{l}}$.\\
b) Quel est le champ magnétique complexe $\overrightarrow{B_{l}}$ associé au champ électrique incident?\\
c) Que vaut, pour l'onde incidente, l'expression du vecteur de Poynting moyen: $\left\langle\overrightarrow{\Pi_{l}(M, t)}\right\rangle=\overrightarrow{\Pi_{0}(M)} ?$

On s'intéresse au champ électrique complexe, en un point infiniment éloigné, associé à un rayon diffracté, dans la direction $\theta_{n}$, telle que $\sin \theta_{n}=n \frac{\lambda}{a}$ avec $n$ entier relatif. On admet que l'amplitude de celui-ci est proportionnelle à $E_{0}$ avec un coefficient de proportionnalité ne dépendant que de $\theta_{n}$ :

$$
\overrightarrow{E_{d}}(M, t)=\underline{\alpha}\left(\theta_{n}\right) E_{0} \overrightarrow{e_{y}} \exp j\left(\omega t-\overrightarrow{k_{d}} \cdot \overrightarrow{O M}\right) .
$$

Q13. a) Comparer les caractéristiques (amplitude, direction et phase) du champ électrique, en un point infiniment éloigné, de l'onde passée par une fente du réseau et de l'onde passée par une autre fente du réseau et diffractée dans la même direction.\\
b) Conclure sur le champ électrique complexe résultant $\overline{E_{d T}}$, c'est-à-dire associé à l'ensemble des rayons passés par les $N$ fentes et diffractés dans la même direction $\theta_{n}$.\\
%%%%%%%%%%%%%%%%%%%%%%%%%%%%%%%%%%%%%%%%%%%
\begin{fig}[label]
\centering
\includegraphics[width=0.5\columnwidth]{2024_11_25_be97bbde3376be1663e6g-10}
\capt{ Dispositif à $N$ fentes.}
\end{fig}
%%%%%%%%%%%%%%%%%%%%%%%%%%%%%%%%%%%%%%%

Q14. a) Que vaut en un point $M$ infiniment éloigné dans la direction $\theta_{n}$ l'expression du vecteur de Poynting moyen $\left\langle\overrightarrow{\Pi_{N}}(M, t)\right\rangle=\overline{\Pi_{N}}(M)$ associé à l'onde résultante issue des $N$ fentes ? On l'exprimera en fonction de $N,|\underline{\alpha}|^{2}$ et de $\overrightarrow{\Pi_{0}}(M)$.\\
b) En déduire le lien entre le rapport H évoqué en Q10.d et le nombre $N$ de fentes éclairées.\\
c) Application : calculer le rapport H pour un réseau $R$ de $p=100$ traits par mm éclairé par un faisceau cylindrique de rayon $r=1,0 \mathrm{~mm}$.\\
d) Indiquer dans quelles directions on obtiendra différents " ordres " de l'onde diffractée en complétant le tableau ci-dessous (en appliquant au réseau $R$ la " formule du réseau " donnée avant la question Q12). On exprimera les résultats en degrés angulaires.

\begin{tabular}{|c|c|c|c|c|}
\hline
Ordre n & 0 & 1 & 2 & 3 \\
\hline
$\theta_{n}$ & $?$ & $?$ & $?$ & $?$ \\
\hline
\end{tabular}

Q15. Exprimer, pour le réseau de pas a, les composantes $k_{n, z}$ et $k_{n, x}$ du vecteur d'onde $\vec{k}$ des directions qui correspondent aux maxima d'intensité, en fonction de $n$, a et de $\lambda$.

Dans le traitement fait précédemment, la largeur $\varepsilon$ des fentes n'a pas été prise en compte. Pour corriger la démarche, on fait l'hypothèse que l'onde émergente du réseau est une superposition d'ondes cohérentes qui, dans le cadre de la diffraction, fournit l'amplitude complexe résultante sous la forme:

$$
\underline{A}(x, z, t)=\alpha_{0} \exp (j \omega t) \sum_{n=-\infty}^{n=+\infty} C_{n} \exp \left(-j\left(k_{n, z} z+k_{n, x} x\right)\right)
$$

avec les coefficients égaux à $C_{n}=\frac{\sin \left(\frac{n \pi x}{a}\right)}{n \pi}$ pour $n$ entier relatif non nul et à $C_{0}=\frac{\pi x}{a}$ pour l'ordre nul et $\alpha_{0}$ une constante réelle .

Q16. Commenter la forme proposée, puis écrire la forme de l'amplitude complexe en utilisant la forme de $\underline{A}(x, z, t)$ donnée ci-dessus et les expressions des composantes du vecteur d'onde obtenues en Q15.

Q17. On suppose que la largeur $\varepsilon$ de chaque fente est petite devant le pas du réseau a.\\
On peut alors montrer que l'amplitude complexe dans le cadre de cette approximation vaut : $\underline{A}(x, z, t)=\alpha_{0} \exp j\left(\omega t-\frac{2 \pi}{\lambda} z\right) \sum_{-q}^{q} C_{n} \exp \left(-j 2 \pi \frac{n}{a}\left(-\frac{\lambda n}{2 a} z+x\right)\right)$ en posant $q$ égal à la partie entière de $\frac{a}{\varepsilon}$.

Dans la suite, on pourra utiliser le fait que puisque $n$ est un entier $\exp j\left(-\pi n^{2}\right)=\exp j(-\pi n)$.\\
a) Écrire l'amplitude $\underline{A}(x, 0, t)$, l'amplitude $\underline{A}\left(x, z=\frac{1}{2} Z_{T}=\frac{a^{2}}{\lambda}, t\right)$ et l'amplitude $\underline{A}\left(x, z=Z_{T}, t\right)$ en fonction de $x, a, n \omega, \lambda$ et de $t$.\\
b) Justifier l'effet Talbot, c'est-à-dire le fait qu'on retrouve l'image du réseau à la distance $z_{T} / 2$ avec un décalage à préciser tandis qu'à la distance $z_{T}$ on retrouve l'image du réseau non décalée (auto-image).

Q18. Commenter la photo 1 et déduire de celle-ci le rapport $\frac{\varepsilon}{a}$ utilisé et vérifier que la longueur Talbot est bien conforme à la formule donnée.

Q19. Problème. On peut établir que la fonction intensité reçue sur un point $M(x, y, D)$ d'un détecteur placé directement à la distance $D$ d'un réseau de période a et dont les $N$ fentes sont de largeur $\varepsilon$ s'écrit:

$$
I(M)=I_{\max }\left(\sin c\left(\frac{\pi \varepsilon x}{\lambda \mathrm{D}}\right)\right)^{2} \cdot\left(\frac{\sin \left(N \frac{\pi a}{\lambda \mathrm{D}} x\right)}{N \sin \left(\frac{\pi a}{\lambda \mathrm{D}} x\right)}\right)^{2}
$$

Commenter ce résultat en faisant apparaître un terme d'interférences et un terme lié à la dimension du motif du réseau.\\
Actuellement, on utilise l'effet Talbot dans le domaine médical comme méthode d'étude des globules rouges permettant un diagnostic concernant les dysmorphies de ceux-ci par observation d'une image en vraie grandeur d'un échantillon sanguin. Les hématies (figure 5 (b)) dans le prélèvement sanguin analysé sont arrangées en couches mono-cellulaires à symétrie hexagonale (figure 5 (a)) formant un réseau permettant d'observer un effet Talbot. L'imagerie permet alors de détecter une maladie potentiellement mortelle le paludisme (parasitose) et des maladies génétiques comme la drépanocytose (globule en forme de faucille) et la micro-sphérocytose (globules sphériques fragiles).\\
%%%%%%%%%%%%%%%%%%%%%%%%%%%%%%%%%%%%%%%%%%%
\begin{fig}[label]
\centering
\includegraphics[width=0.5\columnwidth]{2024_11_25_be97bbde3376be1663e6g-12}
\capt{(a) - Arrangement hexagonal dans une monocouche d'hématies (à gauche).
(b) forme d'une hématie (à droite)}
\end{fig}
%%%%%%%%%%%%%%%%%%%%%%%%%%%%%%%%%%%%%%%

Les globules sont des éléments diffractants de " largeur $2 \mathrm{r}=8 \mu \mathrm{~m}$ ". L'exploitation des images permet de mesurer une valeur de $\mathrm{a}=12 \mu \mathrm{~m}$ en assimilant la période du réseau à la distance entre les centres de deux hématies voisines.

En déduire ce qu'on observerait sur un écran éloigné.\\
À quelle distance faudrait-il se placer pour voir une photo en vrai grandeur (image Talbot) des globules et de leur disposition si on éclaire avec la source monochromatique de longueur d'onde $\lambda=633 \mathrm{~nm}$ ?
}

\bonus{
\subsection{Interférences d'ondes de matière}
Q20. a) Après avoir rappelé la formule de De Broglie, établir la longueur d'onde associée à un faisceau d'électrons accélérés par une différence de potentiel égale à $U=100 \mathrm{~V}$.\\
b) Pourquoi un microscope électronique a-t-il un meilleur pouvoir de résolution qu'un microscope optique conventionnel ?

Après la réalisation d'interférences avec des électrons et des neutrons, des expériences " historiques " d'interférences avec des atomes ont été réalisées en 1991 avec de l'Hélium par Carnel et Mlynek à Constance et en 1992 avec du Néon par Shimizu et Takuma à Tokyo.

Q21. Après avoir cité le théorème d'équipartition de l'énergie, établir la valeur de la vitesse quadratique moyenne u d'un atome de masse $m$ à la température $T$. Indiquer les valeurs manquantes du tableau ci-dessous (cellule notée «? »).

\begin{tabular}{|l|c|c|c|c|}
\hline
Atome & \begin{tabular}{c}
Masse atomique \\
en $\mathrm{g} \cdot \mathrm{mol}^{-1}$ \\
\end{tabular} & \begin{tabular}{c}
Vitesse u \\
à $T_{1}=295 \mathrm{~K}$ \\
\end{tabular} & \begin{tabular}{c}
Vitesse u \\
à $T_{2}=83 \mathrm{~K}$ \\
\end{tabular} & \begin{tabular}{c}
Vitesse u \\
à $T_{3}=2,5 \mathrm{mK}$ \\
\end{tabular} \\
\hline
Hélium & 4 & $?$ & $225 \mathrm{~m} \cdot \mathrm{~s}^{-1}$ &  \\
\hline
Néon & 20 &  &  & $?$ \\
\hline
\end{tabular}

Q22. Dans l'expérience de 1991 (figure 6), la source est un dispositif d'émission thermique d'atomes d'hélium He suivie d'une fente diaphragme $F$. Le montage est celui d'une division du front d'onde de type fentes d'Young. Le plan de la double fente est placé à la distance $L=64 \mathrm{~cm}$ de la fente $F$ et à la distance $L^{\prime}=64 \mathrm{~cm}$ du plan le long duquel est déplacé le détecteur qui comptabilise les impacts de particules.\\
Les atomes d'hélium sont émis soit à la température $T_{1}=295 \mathrm{~K}$, soit à la température $T_{2}=83 \mathrm{~K}$.\\
%%%%%%%%%%%%%%%%%%%%%%%%%%%%%%%%%%%%%%%%%%%
\begin{fig}[label]
\centering
\includegraphics[width=0.5\columnwidth]{2024_11_25_be97bbde3376be1663e6g-13}
\capt{Expérience de 1991}
\end{fig}
%%%%%%%%%%%%%%%%%%%%%%%%%%%%%%%%%%%%%%%

a) On admet qu'il y a diffraction du faisceau de particules au passage par $F$, puis par les fentes $F_{1}$ et $F_{2}$, c'est-à-dire ouverture angulaire du faisceau d'un angle voisin de $2 \lambda / \varepsilon$ avec $\varepsilon$ la largeur de fente. La largeur de la fente $F$ vaut $2 \mu \mathrm{~m}$ et celles des fentes $F_{1}$ et $F_{2}$ valent $1 \mu \mathrm{~m}$. La distance entre elles vaut $\mathrm{a}=8 \mu \mathrm{~m}$. Vérifier que les chemins d'amplitude de probabilité permettent d'envisager deux chemins différents des particules avant leur détection.\\
b) Le détecteur met en évidence des droites parallèles équidistantes sur lesquelles le nombre des impacts est très élevé et entre ces droites, des zones quasiment sans impacts. Interpréter.\\
c) Estimer la distance entre ces droites pour les températures $T_{1}$ et $T_{2}$.

Q23. L'expérience de 1992 est aussi un dispositif interférentiel de fentes d'Young, mais on utilise le champ de pesanteur pour accélérer les atomes de Néon (figure 7).\\
%%%%%%%%%%%%%%%%%%%%%%%%%%%%%%%%%%%%%%%%%%%
\begin{fig}[label]
\centering
\includegraphics[width=0.5\columnwidth]{2024_11_25_be97bbde3376be1663e6g-14}
\capt{Expérience de 1992}
\end{fig}
%%%%%%%%%%%%%%%%%%%%%%%%%%%%%%%%%%%%%%%


Les atomes de Néon, supposés indépendants les uns des autres, sont lâchés sans vitesse initiale et tombent dans le champ de pesanteur $g=9,8 \mathrm{~m} \cdot \mathrm{~s}^{-2}$.\\
a) On mesure un interfrange de $1,1 \mathrm{~mm}$ : que vaut la longueur d'onde associée ? Quelle valeur de la vitesse en déduisez-vous au niveau du détecteur?\\
b) En appliquant les lois de la mécanique classique, évaluer la vitesse des particules en chute libre au niveau des fentes. Que vaut-elle au niveau du détecteur? Commenter.\\
c) Pourquoi a-t-on utilisé des atomes ultra-froids ?
}


\end{exo}
%%%%%%%%%%%%%%%%%%%%%%%%%%%%%%%%%%%
